\section{Finite-dimensional types of finite categories}
\label{sec:D}

This section uses left category modules, that is, covariant functors to
abelian groups.  For a connected category $C$ fix a base object, write
$\pi=\pi_1(BC)$, and choose paths from the base object to every object.
These choices give a functor $\eta:C\to\pi$ to the one-object group
category.  Changing the choices conjugates the corresponding local system
and has no effect on the assertions below.

The general implication from finite-dimensional type to finite CW type
is false. The explicit categories of
Theorems~\ref{D:Q8-final-table} and~\ref{D:C23-final-table} have
nonzero Wall obstruction and fundamental groups $Q_8$ and $C_{23}$,
respectively. Both have executed finite composition tables and CW
homotopy dimension at most eight. The implication is true for finite EI categories,
for fundamental group $C_2$, and for categories with at most two
objects. The last class is classified at the end of the section.

\subsection{The precise finiteness-obstruction reduction}

\begin{proposition}[Finite domination from finite dimension]
\label{D:domination}
If $C$ is a finite category and $BC$ has the homotopy type of a
finite-dimensional CW complex, then $BC$ is finitely dominated.  For each
connected component its Wall obstruction is therefore defined.
\end{proposition}
\begin{proof}
The normalized nerve has finitely many cells in every dimension: an
$r$-simplex is a string of $r$ arrows in a finite set.  Let $Z$ be a
$d$-dimensional CW complex and let $u:Z\to BC$ and $v:BC\to Z$ be
homotopy inverses.  Cellular approximation makes $u$ cellular, so its
image lies in the finite skeleton $K=(BC)^{(d)}$.  Thus $u$ factors as
$Z\xrightarrow{u'}K\xrightarrow{i}BC$, and
$i(u'v)=uv\simeq\operatorname{id}_{BC}$.  This exhibits a domination
by the finite CW complex $K$.  The argument applies componentwise.
\end{proof}

\begin{proposition}[The induced syzygy, with its correct domain]
\label{D:syzygy}
Let $P_*\to\underline{\mathbb Z}$ be the standard projective bar
resolution of the constant left $\mathbb ZC$-module.  Let
$L=\mathbb Z\pi$ be the right $C$-module whose structure maps are
right multiplications by the images under $\eta$, with the commuting
left $\mathbb Z\pi$-action.  Set
$\Omega^n\underline{\mathbb Z}=\operatorname{coker}(P_{n+1}\to P_n)$.
If $BC$ is finitely dominated, then for sufficiently large $n$ the module
\[
 Q_n=L\otimes_{\mathbb ZC}\Omega^n\underline{\mathbb Z}
\]
is a finitely generated projective left $\mathbb Z\pi$-module, and
\[
 w(BC)=(-1)^n[Q_n]\quad\hbox{in }\widetilde K_0(\mathbb Z\pi).
\]
This assertion does not say that
$\Omega^n\underline{\mathbb Z}$ is projective over $\mathbb ZC$.
\end{proposition}
\begin{proof}
The complex $F_*=L\otimes_{\mathbb ZC}P_*$ is the cellular chain
complex of the universal cover, with one free $\mathbb Z\pi$-summand
for each nondegenerate simplex of $NC$.  More explicitly, use the bar
resolution whose summand for a string $x_0\to\cdots\to x_n$ is
$\mathbb ZC(x_n,-)$.  After tensoring, this summand is $L(x_n)$;
all faces retain the lift of the last vertex except the last face,
which applies the right $C$-module structure map of the last arrow.
This is precisely the lifted nerve boundary.  Tensor product is right exact, so
\[
 Q_n=\operatorname{coker}(F_{n+1}\longrightarrow F_n).
\]
Finite domination gives a chain-homotopy equivalence between $F_*$ and
a bounded complex of finitely generated projective
$\mathbb Z\pi$-modules.  To see the relevant truncation explicitly,
choose a chain homotopy $h$ from the identity of $F_*$ to a map factoring
through a bounded complex.  Above its top degree one has
$dh+hd=1$.  In an exact degree this supplies a section
$h:\operatorname{im}d_{n+1}\to F_{n+1}$ of $d_{n+1}$, and, one
degree lower, shows that the cycles there are a direct summand of the
finite free chain module.  Taking $n$ larger than both the homological
and homotopy bounds, $Q_n\cong\operatorname{im}d_n$ is consequently
finitely generated projective.  The truncated projective complex
\[
 0\longrightarrow Q_n\longrightarrow F_{n-1}\longrightarrow
 \cdots\longrightarrow F_0\longrightarrow0
\]
is chain-homotopy equivalent to $F_*$.  Its alternating projective class
is the Wall obstruction; all $F_i$ are free and disappear in reduced
$K_0$.  This proves the formula.
\end{proof}

\begin{proposition}[Why a category-module d\'evissage needs another step]
\label{D:syzygy-counterexample}
There is a four-arrow finite EI category $C$ with contractible $BC$ for
which the constant covariant integral module has infinite projective
dimension.  Thus finite dimensionality of $BC$ cannot justify applying
$K_0(\mathbb ZC)$ to a high category-module syzygy.
\end{proposition}
\begin{proof}
Take objects $a,b$, arrows
$1_a,g:a\to a$, $u:a\to b$, and $1_b:b\to b$.
The nonidentity composition rules are $g^2=1_a$ and $ug=u$.
Associativity follows either from the action of $C_2$ on the singleton
$\{u\}$ or by the complete finite check in
\texttt{code/d\_syzygy\_counterexample.py}.
The object $b$ is terminal, so its identity arrow gives a natural
transformation from the identity functor to the constant functor at
$b$; therefore $BC$ is contractible.

Evaluation at $a$ takes the covariant representable
$\mathbb ZC(a,-)$ to the regular module $\mathbb ZC_2$ and the other
representable $\mathbb ZC(b,-)$ to zero.  It is exact and takes
projectives to projective $\mathbb ZC_2$-modules.  It takes the constant
module to the trivial module $\mathbb Z$.  A finite projective
resolution over $\mathbb ZC$ would therefore give one over
$\mathbb ZC_2$.  But the periodic resolution, with differentials
alternately $g-1$ and $1+g$, gives
$H^{2j}(C_2;\mathbb Z)=\mathbb Z/2$ for every $j\geq1$.
The trivial module cannot have finite projective dimension.

This counterexample does not refute the induced-syzygy formula of
Proposition~\ref{D:syzygy}.  Here $\pi=1$, and the induced high syzygies
are projective over $\mathbb Z$.  What fails is a proposed passage to
the projective class group of $\mathbb ZC$ before induction.
\end{proof}

\subsection{A complete answer for finite EI categories}

\begin{lemma}[A finite groupoid diagram for an EI category]
\label{D:EI-decomposition}
Let $C$ be a finite skeletal EI category.  There are a finite poset $I$,
a functor $G:I\to\mathrm{Cat}$ taking values in finite groupoids, and
a homotopy equivalence
\[
 B\left(\int_I G\right)\longrightarrow BC.
\]
Here $I$ is the poset of nonempty strict chains in the object poset of
$C$, ordered by reverse inclusion, and $G$ is given explicitly below.
\end{lemma}
\begin{proof}
For objects $x,y$ put $x\leq y$ if $C(x,y)\ne\varnothing$.
This is a partial order.  Indeed, if there are arrows both ways, their
two composites are invertible because $C$ is EI; an arrow having both
a left and a right inverse is invertible, and skeletalness then gives
$x=y$.  Denote this finite poset by $P$.

For a strict nonempty chain $\sigma=(x_0<\cdots<x_k)$ let $G_\sigma$
have as objects all strings
$\alpha=(x_0\xrightarrow{\alpha_1}x_1\to\cdots
\xrightarrow{\alpha_k}x_k)$ in $C$.  A morphism
$\alpha\to\beta$ is a tuple
$(h_0,\ldots,h_k)$, $h_i\in\operatorname{Aut}_C(x_i)$, satisfying
$\beta_i h_{i-1}=h_i\alpha_i$ for $1\leq i\leq k$.
All such tuples are invertible and there are finitely many of them.
Deleting vertices, composing the intervening arrows, and deleting
the corresponding entries in the tuple defines a strict functor
$G_\sigma\to G_\tau$ whenever $\tau\subseteq\sigma$.
These functors define $G:I\to\mathrm{Cat}$.

Write $S=\int_I G$.  Its objects are the strings $\alpha$ just
described.  A morphism $\alpha\to\beta$ consists of a face
$\tau\subseteq\sigma$ and an isomorphism
$h:\alpha|_\tau\to\beta$.  Define $p:S\to C$ by taking the first
vertex.  If the first vertex of $\tau$ is $x_j$, define
\[
 p(\tau,h)=h_j\alpha_{0j}:x_0\longrightarrow x_j,
\]
where $\alpha_{0j}=\alpha_j\cdots\alpha_1$ and
$\alpha_{00}=1$.  The commutative-square equations for $h$ give
$p(gf)=p(g)p(f)$; identity tuples give identity arrows.

We prove that $c\downarrow p$ is contractible for every $c\in C$.
Its objects are pairs $(\alpha,\gamma:c\to x_0)$.
If $c<x_0$, prepend the arrow $\gamma$ to $\alpha$.
If $c=x_0$, then $\gamma$ is an automorphism: retain the same
vertices, replace the first arrow $\alpha_1$ by
$\alpha_1\gamma$ when there is a first arrow, and retain the other
arrows.  In either case the new string begins at $c$ and is equipped
with the identity map from $c$.  Call the resulting object $T(\alpha,\gamma)$.

For a morphism in $c\downarrow p$, enlarge both its vertex chains by
$c$ and use the tuple entry $1_c$ at $c$, retaining the given entries
at all larger vertices.  The comma-category equation
$p(f)\gamma=\delta$ verifies the square involving the new first
arrow; all other squares are the original ones.  This defines a
functor $T$.  Deleting a newly prepended vertex, or using the tuple
$(\gamma,1,\ldots,1)$ when $c=x_0$, gives a natural transformation
$T\Rightarrow\operatorname{id}_{c\downarrow p}$.
Deleting every vertex other than $c$ gives a natural transformation
from $T$ to the constant functor at the singleton string $c$ with
structure map $1_c$.  Hence $B(c\downarrow p)$ is contractible.
Quillen's Theorem A \cite[\S1, Theorem A]{quillen1973D} now proves
that $Bp$ is a homotopy equivalence.
\end{proof}

\begin{theorem}[Finite EI categories with finite-dimensional type]
\label{D:EI-finite}
Let $C$ be a finite category in which every endomorphism is invertible.
If $BC$ has the homotopy type of a finite-dimensional CW complex, then
it has the homotopy type of a finite CW complex.  More precisely, the
following construction of a finite acyclic category $K$ depends only
on $C$, and $BC\simeq BK$ under the stated hypothesis:
\[
 K=\int_I\pi_0G,
\]
where $I,G$ are the diagram of Lemma~\ref{D:EI-decomposition} after
passing to a skeleton of $C$.
\end{theorem}
\begin{proof}
Taking a skeleton is a categorical equivalence and preserves
finiteness.  We may therefore use the notation of
Lemma~\ref{D:EI-decomposition}.  Each $\pi_0G_\sigma$ is a finite
set, treated as a discrete category, and the quotient functors give
a natural transformation $G\to\pi_0G$.  The category $K$ is finite.
Its only endomorphisms are identities, and a nonidentity arrow strictly
decreases the length of its underlying chain.  Its nerve thus has
only finitely many nondegenerate simplices, so $BK$ is a finite CW
complex.  In fact $K$ is a poset: between two of its objects there is
at most one morphism.  Its nerve has dimension at most $s-1$, where
$s$ is the number of isomorphism classes of objects of $C$.

Let $X$ be a connected finite-dimensional CW complex.  Miller's
Theorem A \cite[Theorem A, p.~39]{miller1984D} states that
$\operatorname{map}_*(BH,X)$ is weakly
contractible for every locally finite group $H$; in particular this
holds for every finite group.  Evaluation at a basepoint of $BH$ is
a mapping-space fibration with this fibre.  Its section is the map
taking a point of $X$ to the constant map on $BH$.  Consequently the
constant-map inclusion
$X\to\operatorname{map}(BH,X)$ is a weak equivalence.
Every connected component of a finite groupoid is equivalent to a
finite group.  Products over its finitely many components therefore
give a weak equivalence
\[
 \operatorname{map}(\pi_0G_\sigma,X)
 \longrightarrow\operatorname{map}(BG_\sigma,X).
\]
These maps are natural in $\sigma$, since they are induced by the
actual quotient functors, without choices of component representatives.

The classifying-space form of the Grothendieck construction
\cite[Corollary 26.9, derived from Theorem 26.8]{chacholskiScherer2002D} gives
$B\int_I G\simeq\operatorname{hocolim}_I BG$ and the same statement
for $\pi_0G$.  The derived mapping-space property of a homotopy
colimit then identifies mapping spaces into $X$ with the homotopy
limits of the displayed objectwise mapping spaces.  Homotopy limits
preserve objectwise weak equivalences.  Thus the quotient map
$q:B\int_I G\to BK$ induces a weak equivalence of derived mapping
spaces into every such $X$, and in particular a bijection
\[
 q^*:[BK,X]\xrightarrow{\cong}[B\textstyle\int_I G,X].
\]

The map $q$ induces a bijection on components: components of each
groupoid are connected, and each arrow of $K$ lifts after choosing a
representative of its source component.  Thus $BK$ is connected when
$BC$ is connected.  Put $Y=B\int_I G$.  The hypothesis and
Lemma~\ref{D:EI-decomposition} make $Y$ homotopy equivalent to a
finite-dimensional CW complex.  Homotopy invariance of the last
bijection allows $X=Y$.  The identity of $Y$ has a preimage
$r:BK\to Y$, so $rq\simeq1_Y$.  Taking $X=BK$, which is already
finite, gives injectivity of $q^*$.  Since
$(qr)q=q(rq)\simeq q=1_{BK}q$, it follows that $qr\simeq1_{BK}$.
Hence $q$ is a homotopy equivalence, and composing with $Bp$ proves
$BC\simeq BK$.  Disconnected categories are handled one component
at a time.  Miller's 1985 correction \cite{miller1985D} changes neither the statement
nor the hypotheses of his Theorem A.
\end{proof}

\subsection{The case of fundamental group of order two}

\begin{lemma}[Integral projectives over $C_2$]
\label{D:C2-K0}
Every finitely generated projective module over
$R=\mathbb Z[C_2]$ is free.  In particular
$\widetilde K_0(R)=0$.
\end{lemma}
\begin{proof}
Write $C_2=\langle t\mid t^2=1\rangle$.  Evaluation at $t=1$ and
$t=-1$ identifies $R$ with the pullback ring
\[
 R=\{(a,b)\in\mathbb Z\times\mathbb Z:a\equiv b\pmod2\}.
\]
Indeed $u+vt$ maps to $(u+v,u-v)$, whose inverse on pairs of equal
parity is $u=(a+b)/2$, $v=(a-b)/2$.
Let $P=eR^N$ be the image of an idempotent matrix.  Write
$P_+=e_+\mathbb Z^N$, $P_-=e_-\mathbb Z^N$, and
$P_0=e_0\mathbb F_2^N$.  The first two modules are free abelian,
say of ranks $r_+,r_-$.  Their reductions modulo two both identify
with $P_0$, because their embeddings as summands of free modules
split before reduction.  Hence $r_+=r_-=r$.

The equality-of-parity description of $R^N$ and the identity
$e(x_+,x_-)=(e_+x_+,e_-x_-)$ show directly that
\[
 P=\{(x_+,x_-)\in P_+\times P_-:
          \overline{x}_+=\overline{x}_-\text{ in }P_0\}.
\]
Choose integral bases of $P_+$ and $P_-$.  The discrepancy of their
reductions is a matrix $A\in\mathrm{GL}_r(\mathbb F_2)$.
The map $\mathrm{GL}_r(\mathbb Z)\to\mathrm{GL}_r(\mathbb F_2)$
is onto: Gaussian elimination over $\mathbb F_2$ uses row swaps and
row additions, each of which has an invertible integral lift; the
rank-zero and rank-one cases are immediate.  Adjusting one basis
by a lift of $A$ makes the two reductions agree.  The displayed
pullback then becomes the direct sum of $r$ copies of $R$.
\end{proof}

\begin{corollary}[Fundamental group $C_2$]
\label{D:C2-finite}
If $C$ is a connected finite category, $\pi_1(BC)\cong C_2$, and
$BC$ has finite-dimensional CW homotopy type, then $BC$ has finite
CW homotopy type.
\end{corollary}
\begin{proof}
Proposition~\ref{D:domination} makes $BC$ finitely dominated.
Its Wall obstruction lies in the group that
Lemma~\ref{D:C2-K0} proves to be zero.  The vanishing criterion for
Wall's obstruction \cite[Theorem F, p.~66]{wall1965D} gives the
asserted finite CW model.  Its hypotheses are met because
Proposition~\ref{D:domination} supplies a domination by a finite
complex, whose dimension may be taken at least two.
\end{proof}

\subsection{The complete finiteness result with at most two objects}

\begin{theorem}[Two objects force a free fundamental group]
\label{D:two-objects}
Let $C$ be a connected finite category with at most two objects.
If $BC$ has finite-dimensional CW homotopy type, then $\pi_1(BC)$
is a finitely generated free group and $BC$ has finite CW homotopy
type.  More precisely, in the two-object case with no arrows from
$b$ to $a$, put
\[
 A=\operatorname{End}_C(a),\qquad
 B=\operatorname{End}_C(b),\qquad S=C(a,b).
\]
Let $r$ be the number of equivalence classes in the finite set $S$
for the equivalence relation generated by
$s\sim vsu$ with $v\in B$, $u\in A$.  Then
$\pi_1(BC)\cong F_{r-1}$.  If there are arrows both ways, then
$\pi_1(BC)=1$.
\end{theorem}
\begin{proof}
For each object $x$, all loops arising from
$\operatorname{End}_C(x)$ die in $\pi_1(BC,x)$.
Here is the precise input from Part A: the universal group of a finite
monoid is a quotient of a finite subgroup at an idempotent in its
minimal ideal, and its canonical classifying-space map is defined
through the Karoubi envelope.  Miller's Theorem B makes the composite
from that finite subgroup to $BC$ trivial on $\pi_1$ under the
finite-dimensional hypothesis.  Surjectivity onto the universal group
therefore kills every endomorphism loop.  This uses neither a
nonunital inclusion as a monoid homomorphism nor an assertion about
higher homotopy groups.

Suppose first that arrows $f:a\to b$ and $g:b\to a$ exist.
Choose $f$ as a tree edge in the edge-and-triangle presentation of
$\pi_1(BC)$.  Every arrow $h:a\to b$ has associated loop
$f^{-1}h=(gf)^{-1}(gh)$ in the groupoid completion; both factors on
the right are endomorphisms at $a$.  Likewise every arrow in the
other direction and every endomorphism at $b$, after transport along
$f$, is represented by a product of endomorphism loops at $a$.
For example an endomorphism $v$ at $b$ gives
$f^{-1}vf=(gf)^{-1}(gvf)$.  All these loops vanish by the preceding
paragraph.  Thus the whole fundamental group is trivial.

Otherwise, connectedness lets us interchange $a,b$ so that $S$ is
nonempty and $C(b,a)=\varnothing$.  After the endomorphism generators
are killed, a composition relation involving a crossing arrow is
exactly $s=vsu$.  Two crossing arrows cannot be composed, so there
are no other relations.  The groupoid completion is therefore the
free groupoid on two vertices joined by $r$ parallel edges.  Choosing
one of them as a maximal tree identifies its vertex group with
$F_{r-1}$.  The one-object case is the already checked finite-monoid
theorem and gives the trivial group.

By Proposition~\ref{D:domination}, $BC$ is finitely dominated.  Bass,
\emph{Algebraic K-theory} \cite[Chapter IV, Corollary 6.4, p.~213]{bass1968D},
states that for a Dedekind domain $R$ and a free group $F$, every
nonzero finitely generated projective $RF$-module is
$(RF\otimes_R L)\oplus(RF)^{n-1}$ for an invertible $R$-module $L$.
Apply this with the Dedekind domain $\mathbb Z$; every invertible
$\mathbb Z$-module is free of rank one.  Thus every finitely generated
projective $\mathbb ZF_{r-1}$-module is free, and
$\widetilde K_0(\mathbb ZF_{r-1})=0$.
Wall's Theorem F \cite{wall1965D} now turns finite domination into finite homotopy
type.  The group of rank zero is covered by the same argument.
\end{proof}

\begin{corollary}[All two-object EI types of finite dimension]
\label{D:two-EI-types}
The connected finite-dimensional classifying-space types of finite
EI categories with at most two objects are exactly finite wedges of
circles, including the point.  A category with exactly two objects
realizes every such type.
\end{corollary}
\begin{proof}
Theorem~\ref{D:EI-finite} gives a finite poset model whose dimension
is at most one, so its connected realization is a finite graph and
is homotopy equivalent to a finite wedge of circles.  Conversely,
the category with two objects, $r+1$ parallel arrows in one direction,
only identity endomorphisms, and no reverse arrows has nerve a graph
with cycle rank $r$.  All category axioms follow from the fact that
no two nonidentity arrows are composable.
\end{proof}

\subsection{Why a finite filtration alone does not complete the d\'evissage}

\begin{proposition}[An alternating-projector test for proposed d\'evissages]
\label{D:alternating-projector}
Let $R$ be a unital ring, let $E=R^m$, and let
$e:E\to E$ be an idempotent $R$-linear endomorphism.  There is a
nonnegative chain complex $T$ of finite free $R$-modules, with a
three-step filtration by subcomplexes, such that:
\begin{enumerate}
\item every associated graded complex has zero differential and is
obtained from a degreewise finite free integral complex by extending
scalars $\mathbb Z\to R$;
\item $T$ is chain-homotopy equivalent to the projective module
$(1-e)E$ placed in degree zero.
\end{enumerate}
Consequently its reduced projective Euler class is
$[(1-e)E]\in\widetilde K_0(R)$, which need not vanish.  This refutes
a d\'evissage based only on the indicated filtration properties; it
does not produce a finite-category counterexample to D1.
\end{proposition}
\begin{proof}
Put $V_n=E$ for $n\geq0$, with zero differential, and let
$f:V\to V$ have components
$f_n=e$ for even $n$ and $f_n=1-e$ for odd $n$.
Use the homological mapping-cone convention

\[
 K=\operatorname{Cone}(f),\qquad K_n=V_n\oplus V_{n-1},
 \qquad d_K(a,b)=(f_{n-1}b,0).
\]
The summand $V_{-1}$ is zero.  Thus
$H_0(K)=(1-e)E$, whereas, for $n\geq1$,
\[
 H_n(K)=\operatorname{coker}f_n\oplus\ker f_{n-1}
       =(1-f_n)E\oplus f_nE\cong E.
\]
Let $W_0=0$ and $W_n=E$ for $n\geq1$, again with zero
differential.  Define $g:W\to K$ by
\[
 g_n(x)=((1-f_n)x,f_nx)\quad(n\geq1).
\]
It is a chain map because $f_{n-1}f_n=0$, and it induces the
displayed isomorphism $W_n\cong H_n(K)$ in each positive degree.
Set $T=\operatorname{Cone}(g)$.  The mapping-cone homology sequence
gives $H_0(T)=(1-e)E$ and $H_n(T)=0$ for $n>0$.

Explicitly $T_0=E$, $T_1=E\oplus E$, and
$T_n=E\oplus E\oplus E$ for $n\geq2$, with
\[
 d_1(a,b)=eb,\qquad
 d_n(a,b,c)=
 \bigl(f_{n-1}b+(1-f_{n-1})c,\ f_{n-1}c,\ 0\bigr)
 \quad(n\geq2),
\]
omitting the third target summand when $n=2$.  Filter by the first
summand, then the first two summands, then all three.  The differential
strictly lowers this filtration, giving the asserted scalar-associated
graded complexes.  Finally, a bounded-below acyclic complex of
projective modules whose degree-zero homology is projective splits
successively, starting at degree zero.  Applying this elementary
splitting to $T$ gives a chain contraction onto $(1-e)E[0]$.
For an entirely elementary example of nonvanishing, take
$R=\mathbb Z\times\mathbb Z$ and $e=(1,0)$ on $R$; the ranks on
the two factors show that $[(0,\mathbb Z)]$ is nonzero modulo the
class of $R$.  The same construction applies to any specified
idempotent matrix over $\mathbb Z\pi$.
\end{proof}

\subsection{Finite parity models and uniform suspension maps}

\begin{theorem}[Nine-element monoids with prescribed parity of spheres]
\label{D:parity-monoids}
There are explicitly given monoids $M_0,M_1$, each with nine elements,
such that
\[
 BM_0\simeq\bigvee_{k\geq1}S^{2k},\qquad
 BM_1\simeq\bigvee_{k\geq1}S^{2k+1}.
\]
The multiplication and complete integral projective resolutions are
given in the proof.  All multiplication axioms and the repeated
resolution kernels are independently checked by
\texttt{code/d\_parity\_monoids.py}.
\end{theorem}
\begin{proof}
We use algebra order for multiplication: $xy$ means $x$ after $y$.
Start with the six-element monoid
$N=\{1,e_0,e_1,u,v,0\}$.  The element $1$ is its identity and
$0$ its zero.  Regard $e_i$ as the identity at $i$, $u:0\to1$,
and $v:1\to0$.  A product of two of these four arrows is their
composite when composable and at least one is an identity; every
other product is $0$, in particular $uv=vu=u^2=v^2=0$.
This is associative: it is the monoid of paths of length at most
one in the two-arrow quiver, with all longer or noncomposable
products sent to zero and with a global identity adjoined.

Replace $0$ by the rectangular band
$R_0=\{x_{ij}:i,j\in\{0,1\}\}$, whose product is
$x_{ij}x_{kl}=x_{il}$.  Write
$a=x_{00}, b=x_{01}, c=x_{10}, d=x_{11}$.
The underlying set of $M_\epsilon$, $\epsilon\in\{0,1\}$, is
\[
 \{1,e_0,e_1,u,v,a,b,c,d\}.
\]
Products in $N\setminus\{0\}$ retain their value when nonzero
and have value $a$ otherwise.  For
$n\in\{e_0,e_1,u,v\}$ put
\[
 n x_{ij}=
 \begin{cases}x_{ij}&n=e_0,\\x_{0j}&n\ne e_0,\end{cases}
 \qquad
 x_{ij}n=
 \begin{cases}x_{ij}&n=e_\epsilon,\\x_{i0}&n\ne e_\epsilon.
 \end{cases}
\]
The element $1$ acts as identity.  These rules give the whole table.
For a conceptual associativity check, the left action on the row
set is the pointed $N$-action with $e_0$ acting identically and
$e_1,u,v,0$ constantly at $0$; the right action on the column
set has $e_\epsilon$ acting identically and the other nonunits
constantly at $0$.  Both fix the basepoint, their zero actions
are constant, and they commute on separate coordinates.
These facts check every mixed associativity case; the program also
checks all $9^3=729$ triples for each monoid.

Here is a complete resolution over $A=\mathbb ZM_\epsilon$.
In $A$ put
\[
 p=a,\quad q_0=e_0-a,\quad q_1=e_1-a,\quad
 t=1-e_0-e_1+a,
\]
\[
 r=c-a,\quad s=b-a,\quad U=u-a,\quad V=v-a,
 \quad z=d-b-c+a.
\]
The first four elements are orthogonal idempotents summing to $1$.
Together with $r,s,U,V,z$ they form an integral basis of $A$;
the change-of-basis determinant is $\pm1$.  The only nonzero
product of two of the last five basis elements is $rs=z$.
Their idempotent supports are
\[
 r\in q_0Ap,\quad s\in pAq_\epsilon,\quad
 U\in q_1Aq_0,\quad V\in q_0Aq_1,\quad
 z\in q_0Aq_\epsilon.
\]
Thus all calculations below are integral, without a field extension.
The trivial left $A$-module is the rank-one module $S_p$ on which
$p$ acts as $1$ and the other three idempotents act as zero.
Let $S_i$ similarly denote the rank-one module supported at $q_i$.
The projective $Ap$ has basis $(p,r)$; its augmentation to $S_p$
has kernel $\mathbb Zr\cong S_0$.

For $\epsilon=0$, bases for $Aq_0,Aq_1$ are respectively
$(q_0,s,U,z)$ and $(q_1,V)$.  Their radical kernels are
\[
 \ker(Aq_0\to S_0)=A s\oplus\mathbb ZU
                 \cong Ap\oplus S_1,
 \qquad \ker(Aq_1\to S_1)=\mathbb ZV\cong S_0.
\]
The isomorphism $Ap\to As$ is right multiplication by $s$:
it takes $(p,r)$ to $(s,z)$.  Consequently a resolution has
\[
 P_0=Ap,\quad P_{2j+1}=Aq_0\ (j\geq0),\quad
 P_{2j}=Ap\oplus Aq_1\ (j\geq1),
\]
with maps
\[
 d_1(x)=xr,\quad d_{2j}(x,y)=xs+yU,\quad
 d_{2j+1}(x)=(0,xV)\quad(j\geq1).
\]
The displayed kernels identify each image with the next kernel and
prove exactness in every degree.  Tensoring with the trivial right
module sends $Ap$ to $\mathbb Z$ and $Aq_0,Aq_1$ to zero.
All resulting differentials are zero.  This gives one copy of
$\mathbb Z$ in degree zero and in each positive even degree,
and no other homology.

For $\epsilon=1$, bases for $Aq_0,Aq_1$ are
$(q_0,U)$ and $(q_1,V,s,z)$, and
\[
 \ker(Aq_0\to S_0)=\mathbb ZU\cong S_1,\qquad
 \ker(Aq_1\to S_1)=As\oplus\mathbb ZV\cong Ap\oplus S_0.
\]
A resolution is
\[
 P_0=Ap,\quad P_1=Aq_0,\quad P_{2j}=Aq_1\ (j\geq1),
 \quad P_{2j+1}=Ap\oplus Aq_0\ (j\geq1),
\]
where
\[
 d_1(x)=xr,\quad d_2(x)=xU,\quad
 d_{2j+1}(x,y)=xs+yV\ (j\geq1),\quad
 d_{2j}(x)=(0,xU)\ (j\geq2).
\]
Again the listed bases prove exactness indefinitely.  Tensoring
gives one copy of $\mathbb Z$ in degree zero and each odd degree
at least three, with all differentials zero.

To identify the homotopy types, use bar tensors in algebra order.
The following are integral cycles:
\[
 z_{2k}^{(0)}=s\otimes(V\otimes U)^{\otimes(k-1)}\otimes r,
 \qquad
 z_{2k+1}^{(1)}=s\otimes U\otimes
                  (V\otimes U)^{\otimes(k-1)}\otimes r.
\]
Each product of consecutive factors is zero in the augmentation
ideal, and all factors have augmentation zero.
These cycles are primitive generators.  One direct way to check
primitivity is the relative normalized bar resolution for the
subring
$E=\mathbb Zp\oplus\mathbb Zq_0\oplus\mathbb Zq_1\oplus\mathbb Zt$.
Write $A=E\oplus J$, where $J$ has basis $r,s,U,V,z$.
The relative bar resolution has terms
$A\otimes_E J^{\otimes_E n}\otimes_E S_p$; its augmentation
is split over $E$, its usual extra-degeneracy contraction proves
exactness, and its terms are projective left $A$-modules because
every indicated $E$-module has a basis of idempotent-supported
free integral summands.  Tensor with the right module $S_p$.
The coefficient functional of the indicated composable tensor
word is a cocycle: the only nonzero radical product is $rs=z$,
and every term so produced contains $z$, on which this functional
is zero.  It takes value $1$ on the indicated cycle.  The ordinary
bar resolution maps to this relative resolution by quotienting
out the $E$-degeneracies, and the displayed tensors map to the
same tensors.  Hence their ordinary bar homology classes are
primitive, as asserted.

In the universal group of either monoid every idempotent is $1$.
The relations $ue_1=a$ and $ve_0=a$ then also give $u=v=1$.
Thus both classifying spaces are simply connected.  To see that
the primitive cycles are spherical, use the finite categories
$C_n$ constructed in Step~4 of Theorem~\ref{thm:E-wedge}.
Label their adjacent pairs by $(a,c)$, by the reversed sequence
of the relevant interior letters $(a,u)$ or $(a,v)$, and by
$(a,b)$, in chronological order, and label every long arrow by
$a$.  Every required chronological product is $a$, so this is
an actual functor $C_n\to M_\epsilon$.
Its alternating top cycle maps, up to the harmless reversal sign,
to the indicated $z_n^{(\epsilon)}$.
Since $BC_n$ is a wedge of $n$-spheres, this gives a spherical
representative of the generator.  Wedge these representatives
over all the nonzero degrees.  They induce an integral homology
isomorphism between simply connected CW complexes, and hence a
homotopy equivalence.  This proves both claimed types.
\end{proof}

\begin{theorem}[Uniform shifts represented by finite categories]
\label{D:parity-shifts}
There are finite categories $T_0,T_1$, each with three objects and
31 arrows, and functors
\[
 h_0:T_0\longrightarrow M_1,\qquad
 h_1:T_1\longrightarrow M_0,
\]
such that $BT_\epsilon\simeq\Sigma BM_\epsilon$.
Under the homology generators of Theorem~\ref{D:parity-monoids},
$Bh_\epsilon$ sends the suspension of each generator to minus
the next generator.  In particular $Bh_0$ is a homotopy equivalence,
while $Bh_1$ is an isomorphism in every homology degree at least
four and its target has the one additional group $H_2=\mathbb Z$.
\end{theorem}
\begin{proof}
Define monoid homomorphisms $F,G:M_\epsilon\to M_{1-\epsilon}$.
The first fixes $1,e_0,e_1,u,v$ and sends $x_{ij}$ to $x_{i0}$.
The second fixes $1$, sends $e_\epsilon$ to $e_{1-\epsilon}$,
sends the other elements among $e_0,e_1,u,v$ to $a$, and sends
$x_{ij}$ to $x_{0j}$.  Substitution in the multiplication rules
verifies both homomorphism equations; the program checks every
pair as well.  Put $w=u$ if $\epsilon=0$ and $w=v$ otherwise.
For every $m\in M_\epsilon$ one has
\[
 wF(m)=G(m)w,\qquad aF(m)=G(m)a=a.
\]
Thus $w$ and $a$ are two natural transformations $F\Rightarrow G$.

Let $A=\operatorname{im}F$ and $B=\operatorname{im}G$ as
submonoids of the target.  They have respectively seven and four
elements.  The element $a$ is a left zero of $A$ and a right zero
of $B$, so $BA$ and $BB$ are contractible.
Take the Grothendieck construction of the span
\[
 A\xleftarrow{F}M_\epsilon\xrightarrow{G}B.
\]
Its three objects have endomorphism monoids $M_\epsilon,A,B$;
the only crossing hom-sets are $A$ from the first to the second
and $B$ from the first to the third.  Their left actions are
multiplication and their right actions are through $F,G$.
This gives all 31 arrows and their compositions.
The functor $h_\epsilon$ restricts to $F$ on the first object,
to the inclusions on the other two, sends a crossing element
$a'\in A$ to $a'$, and a crossing element $b'\in B$ to $b'w$.
The displayed naturality equation verifies compatibility with
source composition; compatibility with target composition and
identities is immediate.  The program checks all typed pairs
and triples and all these functor equations.

The homotopy-colimit theorem identifies $BT_\epsilon$ with the
double mapping cylinder of $BA\leftarrow BM_\epsilon\to BB$.
Since both endpoints are contractible, this is
$\Sigma BM_\epsilon$ up to homotopy.  The endpoint contractions
using their left and right zero elements give the natural
transformation $a:F\Rightarrow G$; the functor $h_\epsilon$
uses $w$ instead.  Therefore its map on a suspended bar cycle
is the difference of the two prism operators.
For a tensor $(m_1,\ldots,m_n)$ their difference is
\[
 \sum_{i=0}^n(-1)^i
 G(m_1)\otimes\cdots\otimes G(m_i)\otimes(w-a)
 \otimes F(m_{i+1})\otimes\cdots\otimes F(m_n).
\]
Apply this linearly to either displayed primitive cycle.
The term $i=0$ vanishes because $F(s)=0$ in the augmentation
ideal.  Every term with $i\geq2$ vanishes because $G(U)=G(V)=0$,
or, at the last endpoint, because $G(r)=0$.
Exactly the term $i=1$ remains.  Since $G(s)=s$, $F(r)=r$,
and $F$ fixes $U,V$, its value is the negative of the target
primitive cycle, with $U$ prepended when $\epsilon=0$ and $V$
prepended when $\epsilon=1$.  This proves the uniform formula
in every degree.  Both the suspension and the target are simply
connected.  For $\epsilon=0$ the formula is a homology
isomorphism in all degrees, so the homological Whitehead theorem
makes $Bh_0$ a homotopy equivalence.  The assertion for
$\epsilon=1$ follows from the already computed homology.
\end{proof}

\begin{proposition}[Strict parity projectors on an eleven-element model]
\label{D:parity-projectors}
There is an eleven-element monoid $M_{01}$ with
$BM_{01}\simeq\bigvee_{n\geq2}S^n$, equipped with two commuting
idempotent monoid endomorphisms $\rho_0,\rho_1$.  On positive
integral homology they are the projections onto the even and odd
degrees, respectively.  The inclusions $M_0,M_1\to M_{01}$ give
a wedge equivalence represented by a functor from a finite
Grothendieck construction.
\end{proposition}
\begin{proof}
Repeat the construction of Theorem~\ref{D:parity-monoids} with
row set $\{0,1\}$ and column set $\{0,1,2\}$.
The row action is unchanged.  On columns, $e_0$ fixes $1$ and
sends $2$ to $0$, while $e_1$ fixes $2$ and sends $1$ to $0$;
$u,v,0$ send every column to $0$.  All actions fix $0$.
Thus $M_{01}$ has the five elements $1,e_0,e_1,u,v$ together
with the six band elements $x_{ij}$, and the multiplication is
specified by these pointed actions and rectangular-band products.
The same action argument proves associativity; the program
\texttt{code/d\_parity\_projectors.py} prints the full table and
checks all $11^3=1331$ triples.

Use the same orthogonal idempotents $p,q_0,q_1,t$ and elements
$r,U,V$, and put
$s_i=x_{0,i+1}-a$ and $z_i=r s_i$ for $i=0,1$.
Projective bases are
\[
 Ap=(p,r),\qquad Aq_0=(q_0,s_0,U,z_0),\qquad
 Aq_1=(q_1,s_1,V,z_1).
\]
They give, integrally,
\[
 \Omega S_p=S_0,\qquad
 \Omega S_0\cong Ap\oplus S_1,\qquad
 \Omega S_1\cong Ap\oplus S_0.
\]
Consequently the complete resolution has $P_0=Ap$, $P_1=Aq_0$,
and $P_n=Ap\oplus Aq_{(n-1)\bmod2}$ for $n\geq2$.
After tensoring with the trivial right module, its differentials
are zero and there is precisely one integral generator in every
degree at least two.  The two nine-element monoids embed by using
respectively columns $\{0,1\}$ and $\{0,2\}$ and fixing
$1,e_0,e_1,u,v$.  The primitive cycles from the previous theorem
map to the corresponding primitive cycles here.  Primitivity
follows by the same relative-bar coefficient functional.

The finite Grothendieck construction of
$M_0\leftarrow *\to M_1$ has classifying space
$BM_0\vee BM_1$.  Its functor to $M_{01}$ is given by the two
inclusions and the unit on the central object.  It induces an
isomorphism on all integral homology.  Its source and target are
simply connected, the latter by the same idempotent relations
as before, so it is a homotopy equivalence.

Finally $\rho_0$ fixes the nonband elements, fixes column $1$,
and sends column $2$ to column $0$; $\rho_1$ fixes column $2$
and sends column $1$ to column $0$.  These column maps commute
with every pointed $N$-action, hence are monoid homomorphisms.
They are commuting idempotents.  Each retains exactly the primitive
cycles in its designated parity and kills the other cycles
literally in the bar complex.  This proves the assertions.
\end{proof}

\subsection{Explicit nonzero projective classes for the two remaining finite groups}\label{D:explicit-projectives}
\begin{proposition}\label{D:Q8-projector}
Let $R=\Z Q_8$ and $N=\sum_{g\in Q_8}g$. The matrix
\[
 e_8=\begin{pmatrix}1-8N&21N\\-3N&8N\end{pmatrix}
\]
is idempotent, and its image has nonzero class in $\widetilde K_0(R)$.
\end{proposition}
\begin{proof}
The identity $N^2=8N$ verifies all four entries of $e_8^2=e_8$.
The row $(3,-8)$ maps its two columns to $3$ and $-N$, so it maps
$\im e_8$ onto the ideal $(3,N)\subset R$. Both groups have integral
rank eight: on every nontrivial rational representation $e_8$ is
$\operatorname{diag}(1,0)$, and on the trivial representation it is a
rank-one projector. The surjection is therefore an isomorphism.

Here is the precise external nonvanishing input and its applicability.
Martinet's Theorem I.1 \cite[pp.~400--401]{martinet1971} defines the
rank-one projective modules $P_a$ for odd $a$, proves that $P_1=R$,
and proves $P_a\cong P_b$ precisely when $a\equiv\pm b\pmod8$.
His remark on p.~401 identifies $(3,N)$ with $P_3$.
Thus $(3,N)$ is not free. His Theorem II.1, proved on pp.~402--404,
states that every finitely generated stably free $\Z Q_8$-module is
free. These hypotheses apply to the image of our explicit finite
idempotent, so its reduced class cannot vanish. This uses stable
nonfreeness, not merely nonfreeness. Ranicki's explicit projection
calculation \cite[pp.~4--5]{ranickiExplicit} gives the same matrix and
identification independently. Martinet's 2014 correction
\cite{martinet2014correction} concerns the arithmetic sign choice in
Theorem IV.1; it does not alter Theorems I.1 or II.1. We opened the
original proof and the correction.
\end{proof}

\begin{proposition}\label{D:C23-projector}
Let $R=\Z C_{23}=\Z[g]/(g^{23}-1)$, set $N=\sum_{a=0}^{22}g^a$, and put
\[
 \omega=1+g+g^2+g^3+g^4+g^6+g^8+g^9+g^{12}+g^{13}+g^{16}+g^{18}.
\]
Then
\[
 e_{23}=\begin{pmatrix}
 4+8N&\omega-1+N\\
 2\omega-46N&-3-8N
 \end{pmatrix}
\]
is idempotent and $[\im e_{23}]\ne0$ in $\widetilde K_0(R)$.
\end{proposition}
\begin{proof}
Counting sums of the displayed quadratic residues modulo $23$ gives
\[
 N^2=23N,\qquad \omega N=12N,\qquad
 \omega^2-\omega+6=6N.
\]
Substitution proves the idempotent identity entry by entry. Alternatively,
the injective ring map
\[
 R\longrightarrow \Z[\zeta_{23}]\times\Z,
 \qquad g\longmapsto(\zeta_{23},1)
\]
identifies $R$ with the pairs whose residues agree in $\F_{23}$.
The two images of $e_{23}$ are
\[
 \begin{pmatrix}4&\omega-1\\2\omega&-3\end{pmatrix},
 \qquad
 \begin{pmatrix}188&34\\-1034&-187\end{pmatrix}
   =\binom{2}{-11}(94,17).
\]
The second is idempotent since $188-187=1$. In the first ring
$\omega^2-\omega+6=0$; writing
$p=(2,\omega)^{\mathsf T}$ and $q=(2,(\omega-1)/2)$ gives $qp=1$,
so the first matrix is $pq$.

Let $L=\Q(\zeta_{23})$, $K=\Q(\omega)=\Q(\sqrt{-23})$, and
$\mathcal O_K=\Z[\omega]$. The polynomial discriminant is $-23$,
so this is the full ring of integers. The ideal
$I=(2,\omega)\subset\mathcal O_K$ has norm $2$.
It is not principal: a generator would have norm $2$, whereas
\[
 \operatorname{Norm}_{K/\Q}(a+b\omega)
 =a^2+ab+6b^2=(a+b/2)^2+23b^2/4
\]
never equals $2$ for integers $a,b$. The conjugate prime
$(2,\omega-1)$ does not contain $\omega-2$, and this element has
norm $8$. Unique ideal factorization therefore gives
$I^3=(\omega-2)$, so $[I]$ has exact order three.

The extension degree $[L:K]=11$ is prime to three. If
$I\mathcal O_L$ were principal, its ideal norm down to $K$ would
make $I^{11}$ principal, a contradiction. The image of $pq$ is
isomorphic to $(I\mathcal O_L)^{-1}$: the injection
$r\mapsto(2r,\omega r)$ has precisely that fractional ideal as its
integral domain, and $q$ is its inverse on the image because $qp=1$.
A stably free rank-one module over the commutative ring
$\mathcal O_L$ is free by taking its top exterior power. Hence the
extension of $\im e_{23}$ to $\mathcal O_L$ is not stably free,
and neither is $\im e_{23}$ over $R$.
\end{proof}

\begin{remark}\label{D:projector-check}
The program \texttt{d\_explicit\_projectives.py} gives both complete group
tables, verifies all $512$ and $12\,167$ associativity triples, checks
all integral coefficients of both projector identities, and computes
the exact ranks $8$ and $23$ of their integral regular matrices.
It also checks both complementary projectors. Nontriviality of a
projective class is established by the proofs above, not inferred from
a matrix rank calculation. These are algebraic inputs to the finite
category construction; by themselves they are not counterexamples to
Question D1.
\end{remark}


% Finite coefficient realization proof. An executed Q8 table follows in
% sections/D_final_table.tex.

\subsection{A finite realization calculus for the alternating projector}

\begin{lemma}[Finite coefficient data may be tensored with a finite category]
\label{D:finite-coefficient-tensor}
Let $K$ be a finite connected polyhedron.  Consider a finite direct
homotopy-coherent diagram of finite polyhedra under and over $K$,
including only finitely many specified homotopies and higher homotopies.
Here ``over $K$'' is derived: a map $f:Y\to Z$ is accompanied by a
specified homotopy $r_Zf\simeq r_Y$ relative to the section $K$.
Let $A$ be the classifying space of a finite pointed category.
The diagram obtained by tensoring these retractive spaces with $A$,
and then applying finitely many products and finite homotopy colimits,
has a model built from finite categories and functors.  The specified
homotopies have finite categorical representatives in this model.
The assertion also allows a finite list of functors, natural
transformations, and finite Grothendieck constructions in the variable
pointed category $A$.
\end{lemma}
\begin{proof}
We describe all finite constructions needed.  A single map with its
specified over-$K$ homotopy already admits a finite strict model.
If $f:Y\to Z$ fixes $K$ and $H:r_Zf\simeq r_Y$ is the given
homotopy relative to $K$, form its relative mapping cylinder
\[
 Z_f=(Y\times[0,1])\cup_f Z\,/\,((k,t)\sim k\text{ for }k\in K).
\]
Choose the orientation of $H$ so that it is $r_Y$ at the source
end and $r_Zf$ at the target end.  It defines a projection
$Z_f\to K$, agreeing with $r_Z$ on $Z$.  The two inclusions
$Y\to Z_f\leftarrow Z$ are literal maps under and over $K$,
and the second is a weak equivalence on underlying spaces.  It is
therefore invertible in the derived retractive category; the ordinary
cylinder contraction need not preserve the projection to $K$.
Thus they represent the given derived map.  A homotopy between
such data is treated by the same relative mapping-cylinder
construction with an additional interval parameter; a specified
$j$-fold homotopy uses a $j$-cube.  Every space used is finite.
No path-space replacement is being asserted finite.

Here is a relative finite-model construction which avoids any assertion
that an arbitrary finite family of maps admits one simultaneous
triangulation.  First make the finitely many coefficient maps and
parameter homotopies piecewise linear.  Proceed in increasing parameter
dimension: the boundary data have already been made piecewise linear,
the homotopy extension property corrects the old filling to that
boundary, and relative simplicial approximation supplies a piecewise
linear filling.  Only finitely many fillings are used.

Retain a map as the two inclusions into its relative mapping cylinder,
as in the first paragraph, instead of triangulating the map itself.
Treat a specified homotopy by its relative prism and mapping cylinder,
and do the same for the finitely many higher parameter cubes.  Duplicate
each occurrence of an object before this construction.  Attach these
cylinders and prisms in the order of the finite direct pasting shape.
Each old occurrence is a subpolyhedron of the resulting finite
polyhedron: attaching a cylinder or a parameter prism does not identify
distinct points of the old target.  All copies of the specified section
are identified with the one copy of $K$.  The given projection
homotopies agree on the attaching boundaries, so they define one
piecewise linear retraction $r:Z\to K$ on the finite union $Z$.
There are finitely many distinguished subpolyhedra of $Z$, namely the
object occurrences, cylinders, prisms, and their faces.  Their maps are
inclusions.  The target-to-cylinder inclusions which are inverted are
underlying homotopy equivalences.  Thus this finite diagram of
inclusion roofs represents precisely the original finite pasting data,
including the specified homotopies.  This is the cylinder description
of finite bar rectification; finiteness follows also by indexing its
cubes by strings of nonidentity arrows in the finite direct shape.

We record why the single retraction causes no hidden subdivision
cycle.  Start with finite polyhedral decompositions of $Z$ and $K$
on which $r$ is affine, preserving the section and every distinguished
subpolyhedron.  Subdivide $K$ by the images under $r$ of all faces of
the decomposition of $Z$, using the finite arrangement of their
supporting hyperplanes.  Subdivide each old cell of $Z$ by inverse
images of these target cells.  Every resulting cell maps onto one
target cell: its faces are intersections of an old face with the
inverse image of a target face, and the target decomposition was
chosen to subdivide the images of all old faces.  This gives finite
polyhedral face posets $P_Z,P_K$ and an order-preserving map
$P_r:P_Z\to P_K$ taking a cell to its image cell.  On the section
it is the identity.  Every distinguished inclusion commutes strictly
with $P_r$.

These poset models can even be compared over $K$ literally.  Choose
an interior point of each target cell.  In each source cell choose an
interior point over the chosen point of its image cell; surjectivity
onto that cell guarantees such a point.  Use the same choices on all
distinguished subpolyhedra, and on the section use the target choices.
Coning the boundary triangulations to these interior points gives the
usual face-poset homeomorphisms.  Since $r$ is affine on each source
cell, they commute with $r$, and all inclusion comparisons commute
as well.  Consequently the face-poset inclusion roofs are a finite
strict diagram of retractive categories representing the given data.

No self-idempotent on a fixed finite polyhedron has been strictified.
Repeated occurrences are connected by the displayed finite equivalence
roofs.  In particular the retraction triangle is treated as relative
structure on the one map $Z\to K$, not incorrectly declared to be a
direct indexing triangle.

Here is the tensor construction itself.  For a strict retractive
coefficient object $K\xrightarrow{i}Y\xrightarrow{r}K$, $ri=1$,
and a pointed category $C$ with $BC=A$, first take the finite
Grothendieck construction $D$ of
\[
 K\times C\ \longleftarrow\ K\ \longrightarrow\ Y.
\]
The first map uses the base object of $C$.  There are strict functors
$D\to Y\times C$ and $D\to K$, induced respectively by $i$ and
by $r$.  The Grothendieck construction of
\[
 Y\times C\ \longleftarrow\ D\ \longrightarrow\ K
\]
models the derived pushout
\[
 (Y\times A)\ \cup^h_{(K\times A)\cup^h_K Y}\ K.
\]
This is the tensor of $Y$ with $A$ in retractive spaces over $K$.
Every category and every hom-set is finite.  The recipe is strictly
functorial for strict coefficient maps and for functors of pointed
categories.  Products with the finite interval and cubical
coefficient categories carry the specified homotopies along as
well.  The homotopy-colimit theorem transfers the preceding
coefficient equivalences through these constructions.  Additional
finite Grothendieck constructions remain finite, and their
comparison maps and natural transformations are the same
explicit tuple formulas.  This proves the lemma.
\end{proof}

\begin{lemma}[The coefficient matrices and their controlled homotopies]
\label{D:coefficient-matrices}
Let $K$ be a finite connected CW complex of dimension $k$, with
fundamental group $\pi$, and take $m>k+1$, $m\geq3$.
Put $Y=K\vee\bigvee_{j=1}^N S^m$, as a retractive space over $K$.
Every endomorphism of the free $\mathbb Z\pi$-module on the sphere
cells is represented by a map $Y\to Y$ in the derived retractive
homotopy category.  In particular, for $q=1-e$, choose coefficient maps representing
$e$ and $q$ and one derived nullhomotopy for each of $eq=0$ and
$qe=0$.  These two nullhomotopies have finite coefficient
representatives, including their specified over-$K$ homotopies.
Their cellular chain homotopies in degree $m$ can be taken zero.
\end{lemma}
\begin{proof}
The retraction induces a split surjection on every homotopy group.
Its homotopy fibre is simply connected, and the relative Hurewicz
isomorphism for the $m$-cell attachments identifies its $m$-th
homotopy group with the free $\mathbb Z\pi$-module on the attached
sphere cells.  A map under and derived-over $K$ is specified on
these finitely many sphere cells by elements of that group.
Composition acts on their coefficients by matrix multiplication,
with the chosen left/right module convention.  Thus choose maps representing $e,q$ and two paths in the derived
mapping space representing $eq\simeq0$ and $qe\simeq0$.  Each such
path is an ordinary homotopy of the underlying maps accompanied
by its two-parameter over-$K$ coherence.  The source of each
adjoint map is a finite CW complex times finitely many intervals,
so these are finite coefficient data.  We do not assert that an
arbitrary list of equalities in a homotopy category admits
arbitrarily prescribed compatible higher homotopies.

More concretely, write each matrix entry as a finite integral sum
of elements of $\pi$.  Pinch an $m$-sphere into the required number
of summands, use a reflection for a negative coefficient, and
whisker each summand along a finite edge path in $K$ representing
the group element.  The over-$K$ homotopy contracts the path and
its reverse after collapsing the sphere.  The relative co-H
operations are abelian because $m\geq3$; reassociation,
commutation, inverse cancellation and finitely many path
homotopies in $K$ give representatives of the two chosen
nullhomotopies.  No relation between these two nullhomotopies
is required: they occur on different parity summands.
These are finite parameter homotopies, to which
Lemma~\ref{D:finite-coefficient-tensor} applies.

The coefficient complex $Y$ has no cells above dimension $m$.
A cellular homotopy from a sphere map to another sphere map
therefore has zero cellular chain homotopy out of degree $m$.
There are no intervening base cells in degrees $m-1$ or $m$,
by the choice $m>k+1$.  This gives the claimed chain control.
Tensoring such a homotopy with a parity factor preserves the
zero relative chain homotopy, since it acts on the coefficient
factor and is the identity on the parity factor.
\end{proof}

\begin{theorem}[Finite-category realization of projective finiteness obstructions]
\label{D:Wall-realization}
Let $\pi$ be a finitely presented group and let
$e:(\mathbb Z\pi)^N\to(\mathbb Z\pi)^N$ be an idempotent
endomorphism.  There is a finite connected category $\mathcal C_e$
whose fundamental group is $\pi$, whose classifying space is
finitely dominated, and whose Wall obstruction is
\[
 w(B\mathcal C_e)=(-1)^{m+2}[(1-e)(\mathbb Z\pi)^N]
 \quad\text{in }\widetilde K_0(\mathbb Z\pi),
\]
where $m$ may be any sufficiently large chosen integer.
In particular D1 has a negative answer, including for fundamental
groups $Q_8$ and $C_{23}$, by the explicit projectives of
Section~\ref{D:explicit-projectives}.
\end{theorem}
\begin{proof}
Choose a finite connected presentation complex $K$ with group $\pi$
and choose $m$ as in Lemma~\ref{D:coefficient-matrices}.
Use the split parity space
$A=BM_0\vee BM_1$, with the finite categorical wedge model of
Theorem~\ref{D:parity-monoids}.  Tensor the coefficient object
$Y=K\vee\bigvee^N S^m$ with $A$ in retractive spaces over $K$.
Its underlying space has type
\[
 V=K\vee\bigvee_{j=1}^N(S^m\wedge A).
\]
Its relative universal-cover chains are homotopy equivalent to
one free $(\mathbb Z\pi)^N$ in every degree at least
$d=m+2$, with zero differential.  Choose the basis supplied by
the primitive spherical generators of the parity monoids.

Let $f:V\to V$ have coefficient $e$ on the $M_0$ summands and
$q=1-e$ on the $M_1$ summands.  Let $f^c$ use $q$ on the first
parity and $e$ on the second.  These are built from the same two
chosen coefficient maps, with their positions exchanged.  It is a map in the derived retractive
category, represented by finite coefficient data tensored with
the parity categories.  The uniform shift maps of
Theorem~\ref{D:parity-shifts} give
\[
 \sigma:\Sigma_K V\longrightarrow V.
\]
They exchange the parity summands, sending each suspended chosen
generator to a unit $\pm1$ times the next generator. No common
suspension-sign convention is needed below. The shifts act on the
parity factor; the coefficient maps act on $Y$.

Form the relative mapping cone
$C=\operatorname{Cof}_K(f)$, and let $j:V\to C$ denote its
inclusion.  The finite coefficient nullhomotopy
$ff^c\simeq0$ defines a map
\[
 \ell:\Sigma_K V\longrightarrow C
\]
whose composite with the cone quotient is $\Sigma_Kf^c$.
One can construct it directly on the two suspension cones:
on one cone use $f^c$ into the source cone of $C$, and on the
other use the specified nullhomotopy.  Thus no stable exactness
assertion is needed to produce this unstable map.
Put
\[
 g=\ell+jf^c\sigma:\Sigma_KV\longrightarrow C,
 \qquad T=\operatorname{Cof}_K(g).
\]
The sum is the relative co-H sum. The explicit table below instead
uses a pinch with branch degrees $(1,-1)$; this only changes a unit
sign in the following calculation. Every construction here is
a finite product, finite homotopy colimit, finite coefficient
homotopy, or one of the explicitly represented parity shifts.
Lemma~\ref{D:finite-coefficient-tensor} therefore supplies a
finite category $\mathcal C_e$ and a homotopy equivalence
$B\mathcal C_e\simeq T$.

We calculate the chains, retaining the specified homotopy.
Write $E=(\mathbb Z\pi)^N$ and renumber the bottom sphere degree
$d$ as zero.  On the relative formal chain model the components
of $f$ alternate $e,1-e$.  The cone $C/K$ has the mapping-cone
complex $K$ of Proposition~\ref{D:alternating-projector}.
The lower coordinate of $\ell$ in degree $n>0$ is
$1-f_{n-1}=f_n$.  Its upper chain-homotopy coordinate is zero
by the finite coefficient control in
Lemma~\ref{D:coefficient-matrices}. The summand $jf^c\sigma$ has
upper coordinate $\epsilon_n(1-f_n)$ and lower coordinate zero,
where $\epsilon_n\in\{1,-1\}$. Hence $g$ has the form
\[
 x\longmapsto(\epsilon_n(1-f_n)x,f_nx).
\]
This is an isomorphism from $E$ to
$(1-f_n)E\oplus f_nE=H_n(C/K)$, with inverse
$(u,v)\mapsto\epsilon_nu+v$. Thus the proof of
Proposition~\ref{D:alternating-projector}, with this signed
coordinate, shows that $C_*(\widetilde T,\widetilde K)$ is
chain-homotopy equivalent to $(1-e)E[d]$.

All spaces in the construction have the same fundamental group
as $K$: their newly attached factors are simply connected, and
the two relative cones apply van Kampen to maps inducing the
identity on $\pi$.  Thus $\pi_1(T)=\pi$.
The universal-cover complex of $K$ is finite free, while the
relative complex just computed is perfect.  The exact triangle
of the pair therefore makes $C_*(\widetilde T)$ chain-homotopy
equivalent to a bounded complex of finitely generated projectives.
Together with the finite presentation of $\pi$, Wall's finite-
domination criterion \cite{wall1965D} makes $T$ finitely dominated. The
additivity of the projective Euler class gives exactly
$(-1)^d[(1-e)E]$, since the finite base contributes zero in
reduced $K_0$.  A finitely dominated CW complex has a finite-
dimensional CW model, for example the mapping telescope of a
homotopy idempotent on a finite dominating complex.  This is the
required finiteness property.
\end{proof}


\subsection{An executed finite composition table with nonzero Wall class}
\label{D:final-table}

The following table is specified by integer indices and total lookup
formulas. Its size makes a rectangular printed array impractical; the
index formulas specify every entry, including every typed composite.
There is no unspecified triangulation or search for a subdivision in
this definition.

\begin{theorem}[An explicit negative answer to D1]
\label{D:Q8-final-table}
There is a connected finite category $\mathcal W_8$ with
\[
\begin{split}
 |\operatorname{Ob}\mathcal W_8|
   &=489\,754\,140\,643\,145\,442,\\
 |\operatorname{Mor}\mathcal W_8|
   &=337\,126\,384\,361\,511\,151\,079\,802.
\end{split}
\]
Its classifying space is finitely dominated and has a CW model of
dimension at most eight. Its fundamental group is $Q_8$, and
\[
 w(B\mathcal W_8)=[\operatorname{im}e_8]\ne0
 \quad\text{in }\widetilde K_0(\mathbb ZQ_8).
\]
Consequently $B\mathcal W_8$ is not homotopy equivalent to a finite
CW complex. The complete indexed table is implemented in
\path{code/d_final_wall_category.py}.
The function \texttt{build\_final\_category()} returns it as
\texttt{Final}; its four operations are
\texttt{source}, \texttt{target}, \texttt{identity}, and
\texttt{compose}.
\end{theorem}

\begin{proof}
We give the finite inputs, all table constructors, and the homotopy and
chain calculations. The implementation is a literal indexing of these
formulas, rather than a criterion used to infer a homotopy type from
a finite homology range.

\paragraph{The group, base, and projector.}
Number $Q_8$ as $2a+s$, where $a=0,1,2,3$ denotes $1,i,j,k$ and
$s=0,1$ is the sign. Multiply signs modulo two; equal nonzero axes
have product $-1$, and $(i,j),(j,k),(k,i)$ have respectively products
$k,i,j$, with the reversed products negative. This specifies all
64 products. Let $L$ be the boundary of the four-dimensional cross
polytope with these eight vertices, the opposite pairs being
$\{g,-g\}$. Left multiplication acts freely on $L$.
Choose the lexicographically least face in each orbit. The category
$K_8$ has these ten representatives as objects and arrows
\[
 (F,H,g):F\longrightarrow H\qquad(gF\subseteq H).
\]
Its identity is $(F,F,1)$ and its composition is
\[
 (H,J,h)(F,H,g)=(F,J,hg).
\]
There are 68 arrows. Its nerve is the orbit of the face-poset nerve
of $L$; hence $BK_8\cong|\operatorname{sd}L|/Q_8$ and has
fundamental group $Q_8$ and universal cover $S^3$. Its nondegenerate
simplex counts are $(10,58,96,48)$ in dimensions zero through three.

Put $R=\mathbb ZQ_8$, $N=\sum_g g$, and use the unimodular change
of basis
\[
 P=\begin{pmatrix}2&5\\1&2\end{pmatrix},\qquad
 e=P^{-1}e_8P=
 \begin{pmatrix}-4&N-10\\N+2&5\end{pmatrix},\qquad q=1-e.
\]
Direct multiplication gives
$eq=qe=(8N-N^2)I=0$ and $e^2=e$. Thus
$[\operatorname{im}e]=[\operatorname{im}e_8]\ne0$ by
Proposition~\ref{D:Q8-projector}.
Each of $e,q$ is a sum of 35 signed matrix units with group labels.

\paragraph{Finite coefficient sets and their exact maps.}
Let $J_n$ have vertices $v_0,\ldots,v_n$ and edges
$u_1,\ldots,u_n$, with the two nonidentity arrows
$v_{i-1}\to u_i\leftarrow v_i$. Put
$A_n=NJ_n/\{v_0,v_n\}$. The quotient identifies the two endpoint
constant simplicial subsets with the basepoint. In degree $k$ it has
$2n(k+1)$ simplices, including the basepoint; monotone pullback acts
on its vertex strings. Take $S=A_1^{\wedge5}$ and
$Y=|NK_8|\vee S\vee S$, using the actual finite simplicial nerve
of $K_8$ as the base and literal pointed simplicial quotients for
the spheres. All spaces below are finite simplicial sets with this
base as a simplicial subset.

A group label $g$ uses the vertex path $(1,g)$, except that $1$ uses
$(1,1)$ and $-1$ uses $(1,i,-1)$. Its interval face category maps
to $K_8$ by face-orbit normalization. A label map from sphere $j$
to sphere $i$ is represented by the cospan whose $i$-th target sphere
has its attachment point transported along this interval. The left
leg uses its initial endpoint, the right leg its final endpoint, and
all unused source spheres collapse to the base. The right leg is a
weak equivalence since the transport interval is contractible.
A negative sign reflects suspension coordinate four, numbered from
zero. All additive pinches initially use coordinate zero, so signed
elementary maps commute strictly with these pinches.

The full endpoint functors, including their projections to $K_8$,
are given by the following finite rule. For a path or disk face $F$
with labelled vertices let $\bar F$ be its least orbit representative,
and choose the least $t_F\in Q_8$ satisfying $t_F\bar F=F$.
Send an inclusion $F\subseteq H$ to
$(\bar F,\bar H,t_H^{-1}t_F)$. Identity and composition follow by
cancellation of the adjacent $t$'s. Each multiplication disk has
boundary the concatenated paths for $h$, $g$, and $(hg)^{-1}$,
and is coned to the least cross-polytope vertex whose opposite is
absent from that boundary. The resulting face-poset disk and its
projection use the same rule. There are exactly 64 fixed disks.
The inner label is used first: an inner $h$ followed by an outer
$g$ uses the disk $(h,g)$.

We spell out the finite homotopy operations combining these data.
The map $A_n\to A_1$ retaining only its last edge has degree one,
while collapsing all subdivision vertices gives
$A_n\to\bigvee_n A_1$. Together they give the nary pinch roof.
For retained-edge collapses $q_a,q_b:J_n\to J_1$, assume $a<b$.
There is a common upper functor $h_{ab}$: vertices below $a$ go
to the left endpoint, vertices from $a$ through $b-1$ to the
edge, and vertices at least $b$ to the right endpoint; edges from
$a$ through $b$ go to the edge, and the remaining edges to the
corresponding endpoint. Thus
$q_a\Rightarrow h_{ab}\Leftarrow q_b$, fixing both endpoints.
Apply these functors independently to branches of an
$n$-by-$n$ suspension grid. Moving all retained edges at once,
reinterpreting the resulting sparse matching by columns, and
reversing the procedure gives any prescribed permutation in four
prism stages. The same construction compares addition in
coordinates zero and four in two prism stages.

Successive pinches of arities $n,m$ are compared to the flat
$nm$ pinch by retaining edges $m,2m,\ldots,nm$ and splitting
$J_{nm}$ into its $n$ consecutive blocks of $m$ edges. These maps
satisfy literally
\[
 q_ns=q_{nm},\qquad p_ns=(\bigvee q_m)t,\qquad
 (\bigvee p_m)t=p_{nm}.
\]
A three-edge mapping cylinder realizes the comparison without
inverting a simplicial map. Finally adjacent opposite terms fold
their two edges into one based $J_1$; the pointed contraction
$\operatorname{id}\Rightarrow\mathrm{peak}\Leftarrow\mathrm{constant}$
contracts every pair simultaneously. A zero tail is sent directly
to the common basepoint. Different coefficient apices are kept in
a target-star mapping cylinder, so this operation never assumes a
strict inverse to a weak equivalence.

All mapping cylinders have explicit simplices: besides endpoint
simplices, a nonconstant cylinder simplex is a pair $(x,c)$ with
$1\leq c\leq k$, representing the monotone map $[k]\to[1]$ whose
first $c$ values are zero. Pullback counts how many selected
vertices lie before $c$; zero or $k+1$ selected zeros use the
specified endpoint map. Relative cylinders identify the base
cylinder with the base. This gives an exact finite normal form
and its pullback function. A chain of cylinders uses an additional
finite factor index. Products, pointed quotients, and wedges act
coordinatewise, retaining the summand tag unless the simplex
hits the base. These are the complete simplicial lookup rules
in the coefficient compiler.

Expanding the two 35-term matrices gives 1225 elementary products
for each order. A mismatched matrix index gives zero; a matched
index uses its fixed group disk. Sorting by matrix entry and group
label, with positive and negative terms adjacent, gives 294 inverse
pairs and 637 zero terms for $eq$, and 300 inverse pairs and 625
zeros for $qe$. The preceding fixed operations give exactly 50
coefficient roofs and 49 actual weak comparison maps for each
nullhomotopy. Every comparison preserves both true $Y$ endpoint
maps. All of this finite input has been executed, including the
two orders and all group/sign/index combinations.

Take the relative mapping-cylinder homotopy colimit $Z_{eq}$ or
$Z_{qe}$ of the corresponding 50-term zigzag. It contains its
initial composite apex and is weakly equivalent to $Y$. The left
legs give an actual map
\[
 H:(Y\times I)/((k,t)\sim k)\longrightarrow Z,\qquad
 H|_0=f'c',\quad H|_1=i_Zr_Y.
\]
Here $I$ is the finite oriented path supplied by the zigzag.
The base collapse is justified: before collapse $K\times I$ is
a simplicial subcomplex of the ordinary mapping cylinder.
Its projection to $K$ is a homotopy equivalence, and pushout along
that subcomplex inclusion preserves the equivalence.
Both endpoint inclusions of the relative domain are weak
equivalences. All coefficient apices have dimension at most
seven; this last path cylinder has dimension at most eight.

\paragraph{The finite category table constructors.}
For each coefficient simplicial set $X$ use the opposite truncated
simplex category
\[
 O_8(X)=(\Delta_{\leq8}\mathbin{\downarrow}X)^{\rm op}.
\]
An object is $(k,x)$, $0\leq k\leq8$, $x\in X_k$. An arrow
from $(k,x)$ to $(l,y)$ is a monotone map
$\theta:[l]\to[k]$ with $y=\theta^*x$. Identity is the identity
ordinal map; composition is composition of ordinal maps. Its
object and arrow counts are respectively
\[
 \sum_{k=0}^8|X_k|,\qquad
 \sum_{k,l=0}^8|X_k|\binom{k+l+1}{l+1}.
\]
Order objects by degree and then by the simplex normal forms;
order outgoing arrows by target degree and the lexicographic
ordinal map. These rules give unique integer indices. The
truncated-simplex comparison of Lemma~\ref{lem:B-prelim} gives
$BO_8(X)\simeq|X|$. To recall the relevant proof, in the comma
category from a simplex to the truncated simplex category,
ordered-image factorization retracts to the surjective ordinal
maps; the unique nondegenerate ancestor is terminal there and
has dimension at most eight. Quillen's Theorem~A compares the
truncated and full simplex categories. The full simplicial
replacement has the homotopy type of $|X|$. Passing to the
opposite category preserves realization. Our convention makes
the number of arrows out of an object depend only on its degree.

Products use ordered pairs and coordinatewise composition.
For a span $A\xleftarrow{F}D\xrightarrow{G}B$, the Grothendieck
category has the three tagged object sets $D,A,B$. Besides
internal arrows it has crossing arrows $(x,a)$ from $x\in D$
to $a_1\in A$ for $a:F(x)\to a_1$, and analogous $B$ arrows.
Composition after a crossing is target composition; composition
before a crossing replaces $(x,a)$ and $u:x'\to x$ by
$(x',aF(u))$. There are no two successive crossings.
The same formulas apply to the four crossings of the two
five-stage pinch diagrams. Identity and associativity therefore
follow from the input laws and functor equations. This is a
universal proof of the table laws, independent of numerical checks.

For exact indexing, partition objects into the finite types used
by these constructors. If a type contains $n_t$ objects and has
$a_t$ outgoing arrows per object, enumerate its objects
consecutively and each object's outgoing arrows in the displayed
internal/crossing order. The block has $n_ta_t$ arrows.
Product types are pairs; span types add their stage tag.
On a crossing, the target arrow is enumerated in the outgoing
block of the image of the source under the given functor.
Decoding these finite mixed-radix indices and applying the
composition rule gives a total lookup for every typed pair.
The table for $\mathcal W_8$ has 3897 such types.

\paragraph{The final finite construction.}
Apply the relative tensor of Lemma~\ref{D:finite-coefficient-tensor}
to the coefficient objects with the two nine-element monoids
$M_0,M_1$. Its two finite Grothendieck spans are strictly
functorial. Form the relative wedge $V$ of the two $Y$ tensors.
Let $W$ be the wedge of the $q$-apex tensor in parity zero and
the $e$-apex tensor in parity one, and let $Z$ be the wedge of
the $Z_{eq},Z_{qe}$ tensors. The initial three factors of each
composite coefficient chain give actual functors
\[
 c':V\longrightarrow W,\qquad f':W\longrightarrow Z.
\]
There are weak right-leg functors $V\to W$ and $V\to Z$.
Under them $c'$ has coefficients $q,e$ and $f'$ has coefficients
$e,q$. The latter assertion uses the explicit junction cylinder
between the two coefficient blocks; its two endpoint maps are
not asserted literally equal.

Tensor the actual $H$ above and form the relative cones
$C=\operatorname{Cof}_K(f')$, $C_{fc}$, and $C_H$.
The endpoint inclusions give equivalences
\[
 C_{fc}\longrightarrow C_H\longleftarrow C_{\mathrm{tensor}(0)}.
\]
The tensor of a coefficient zero is not literally the section
followed by retraction. It factors through $T_K(M_\epsilon)$.
In that base tensor a finite middle functor moves the outer
product stage into the central boundary stage, and has natural
transformations to both the identity and the section--retraction
functor. On the boundary span the middle functor fixes the
product stage and sends its two copies of $K$ to the central
copy. The transformation components are its canonical span
arrows. Both transformations are over $K$ and identity on the
chosen section. Composing with the coefficient retraction and
section gives the further cone equivalences
\[
 C_{\mathrm{tensor}(0)}\longrightarrow C_{\mathrm{middle}}
 \longleftarrow C_0.
\]
For $a\Rightarrow b$, the cone comparison runs
$\operatorname{Cone}(b)\to\operatorname{Cone}(a)$; on a crossing
it sends $t:b(x)\to y$ to $t\eta_x:a(x)\to y$.
These are the directions used here.

The suspension-summand inclusion $\Sigma_KV\to C_0$, these
comparisons, and the strict square $(c',1)$ from $C_{fc}$ to $C$
give a finite cospan representing $\ell:\Sigma_KV\to C$.
A second represents $\psi=jf^c\sigma$. Here $\sigma$ uses the
31-arrow parity shifts of Theorem~\ref{D:parity-shifts}, the
natural finite suspension--tensor comparison, and the relative
wedge; $f^c$ is the complementary coefficient roof.
The implementation retains all base paths. Composition uses
a span homotopy colimit, appending its two canonical natural
transformations to the base path.

The intrinsic pinch uses the height-one diagrams
\[
 K\longleftarrow V\longrightarrow V\longleftarrow V
       \longrightarrow K,\qquad
 K\longleftarrow V\longrightarrow K\longleftarrow V
       \longrightarrow K.
\]
In the first the two middle arrows are identities; all remaining
arrows are retractions. Collapse its first half to obtain a weak
map to $\Sigma_KV$, and retract its middle $V$ to obtain a map
to the second diagram. The ordinary wedge of two suspensions
maps weakly to the second diagram, reversing the ends of the
second suspension. This gives a finite pinch cospan with branch
degrees $(1,-1)$. The wedge/fold construction, with collars on
the stored base paths, gives
\[
 \Sigma_KV\xrightarrow{g_l}G\xleftarrow[\simeq]{g_r}C
\]
representing $g=\ell-\psi$ in this convention.
Every map is strictly over $K$; base paths are retained when a
section is preserved only up to homotopy.
Define $\mathcal W_8$ as the finite Grothendieck construction
of $G\xleftarrow{g_l}\Sigma_KV\to K$. Since $g_r$ is a weak
equivalence over $K$, its realization is the required relative
homotopy cofiber of $g$.

\paragraph{Universal-cover chains and the obstruction.}
The relative universal-cover complex of $V$ is chain-homotopy
equivalent to one free $R^2$ in every degree $n\geq7$, with
zero differential. The coefficient $S^5$ supplies the shift
by five; the parity generators start in degrees two and three.
On this model $f_n$ alternates $e,q$, while its complement
alternates $q,e$. The coefficient nullhomotopies have zero
degree-six cellular chain homotopy after transfer to
$K\vee2S^5$, which has no six-cells. Tensoring with the unchanged
parity factor preserves this zero homotopy.

The additional tensor-zero contraction has the same property
for a functorial reason. It is obtained by first applying the
coefficient retraction $Y\to K$, then the contraction in
$T_K(M)$, and then the coefficient section. On relative
coefficient chains the first map is the zero map from the two
five-sphere cells. The relative tensor chain functor is the
tensor product of this coefficient complex and the reduced
parity complex; hence the entire resulting chain homotopy is
zero on that model. The contraction's naturality, rather than
merely the acyclicity of its intermediate target, justifies
this assertion. Base collars factor through the base and
likewise contribute nothing to the relative complex.

Renumber degree seven as zero and put $E=R^2$. The first cone
has $K_n=E\oplus E$ for $n>0$, $K_0=E$, and
$d(a,b)=(f_{n-1}b,0)$. The lower coordinate of $\ell_n$ is
$1-f_{n-1}=f_n$, and its upper coordinate is zero.
The second summand has upper coordinate
$\epsilon_n(1-f_n)$ and lower coordinate zero, where
$\epsilon_n=\pm1$ includes the parity-shift and pinch
orientations. Therefore $g$ induces the isomorphism
\[
 E\longrightarrow H_n(K)=(1-f_n)E\oplus f_nE,\qquad
 x\longmapsto(\epsilon_n(1-f_n)x,f_nx)\quad(n>0),
\]
with inverse $(u,v)\mapsto\epsilon_nu+v$.
Here is an explicit chain deformation retraction, so this is
stronger than a homology calculation. Write $p_n=f_n$ and
$q_n=1-f_n$, so $p_{n-1}=q_n$. Project the first cone onto its
zero-differential homology summand by
$\pi_n(a,b)=(q_na,p_nb)$, with $H_0=q_0E$.
The evident inclusion $\iota$ and the degree-one map
\[
 s_n(a,b)=(0,p_na)\in K_{n+1}
\]
satisfy $ds+sd=1-\iota\pi$, including the bottom degree.
The attaching map lands in this chosen homology summand.
Its cone therefore splits into contractible two-term pairs in
every positive degree and leaves exactly $(1-e)E$ at the bottom.
Restoring the shift gives the claimed contraction onto
$(1-e)E[7]$. Independent signs on both coordinates would have
the same proof, using inverse $(u,v)\mapsto\alpha_nu+\beta_nv$.

All relative pieces are simply connected above the base, and
the two cones apply van Kampen to maps inducing the identity
on $\pi_1K$. Thus $\pi_1(B\mathcal W_8)=Q_8$.
The finite base has universal-cover free chain ranks
$(10,58,96,48)$ in degrees zero through three. The connecting
chain map from the relative projective module in degree seven
to its suspension vanishes for dimension reasons. Consequently
\[
 C_*(\widetilde{B\mathcal W_8})\simeq
 C_*(\widetilde{BK_8})\oplus(1-e)R^2[7]
\]
over $R$. For completeness there is a direct domination argument.
The nerve has finitely many cells in each degree, so its
seven-skeleton $L_7$ is a finite CW complex. The inclusion into
$B\mathcal W_8$ is seven-connected. Replace it by a fibration;
its fibre is six-connected. Every obstruction to a homotopy
section lies in $H^j(B\mathcal W_8;\mathcal A)$ for a local
coefficient system $\mathcal A$ and $j\geq8$.
These groups all vanish by the bounded projective chain model
just computed. Inductive obstruction theory gives a homotopy
section, so $L_7$ dominates $B\mathcal W_8$.
The mapping telescope of the resulting homotopy idempotent
on $L_7$ is a CW model of dimension at most eight.
This also verifies directly the hypotheses in Wall's
finite-domination criterion \cite{wall1965D}.
Its reduced projective Euler class is
\[
 (-1)^7[(1-e)R^2]=[eR^2]=[\operatorname{im}e_8]\ne0.
\]
Wall's obstruction theorem \cite[Theorem F]{wall1965D}
excludes a finite CW model.

\paragraph{Verification and exact size.}
The ordinal, product, span, and height-one composition rules
prove all category laws universally. Complete small tables
independently check those formulas. The nary flattening audit
checks 6\,765\,548 simplex/pullback equations, 3424 endpoint
equations, and 12\,280 inverse equations in a nontrivial small
instance. The actual $Q_8$ coefficient paths check all 64 group
disks and 512 group/sign/index cases. The final run checks 558
base-retraction equations, 468 base-transformation equations,
and 15\,588 identity/associativity equations at the two integer
ends of every final type block. The last checks sample the
enormous table; they are not an exhaustive enumeration.
The exact type-count recursion gives the counts in the theorem,
and also
\[
\begin{split}
 \#\{\text{typed pairs}\}&=
 138\,512\,978\,818\,954\,605\,575\,263\,121\,838,\\
 \#\{\text{typed triples}\}&=
 52\,322\,049\,455\,781\,789\,059\,845\,756\,837\,400\,964.
\end{split}
\]
The machine-readable record is
\texttt{outputs/d\_final\_wall\_category\_summary.json}.
\end{proof}

\subsection{A second executed table, with fundamental group $C_{23}$}
\label{D:C23-final-table-section}

\begin{theorem}\label{D:C23-final-table}
There is a connected finite category $\mathcal W_{23}$ with
\[
\begin{split}
 |\operatorname{Ob}\mathcal W_{23}|
   &=1\,021\,805\,843\,373\,670\,056,\\
 |\operatorname{Mor}\mathcal W_{23}|
   &=703\,363\,835\,733\,950\,740\,370\,112.
\end{split}
\]
Its fundamental group is $C_{23}$. Its classifying space is finitely
dominated, admits a CW model of dimension at most eight, and has
nonzero Wall finiteness obstruction. In particular the $C_{23}$ case
of D2 is negative, with a fixed finite composition table.
The constructor is \path{code/d_c23_wall_category.py}:
\texttt{build\_c23\_category()['Final']} returns the table with
the operations \texttt{source}, \texttt{target}, \texttt{identity},
and \texttt{compose}. Its object and arrow sets are the consecutive
integer intervals with the displayed cardinalities.
\end{theorem}

\begin{proof}
We specify the new finite inputs and the arithmetic indexing, then
apply the fully typed construction and signed chain calculation in
Theorem~\ref{D:Q8-final-table}. This instance was executed separately;
its definition does not alter the quaternion table.

\paragraph{The shorter integral lift.}
Set $R=\Z C_{23}$ and use the elements $N,\omega$ of
Proposition~\ref{D:C23-projector}. Define
\[
 e=\begin{pmatrix}
 4&\omega-1-N\\
 2\omega-N&-3
 \end{pmatrix},\qquad q=1-e.
\]
The three exact group-ring identities
\[
 N^2=23N,\qquad N\omega=12N,\qquad
 \omega^2-\omega+6=6N
\]
give
\[
\begin{split}
 (\omega-1-N)(2\omega-N)
 &=2\omega^2-2\omega-3N\omega+N+N^2\\
 &=-12.
\end{split}
\]
Thus $\operatorname{tr}e=1$, $\det e=0$, and the two-by-two
Cayley--Hamilton identity over the commutative ring $R$ gives
$e^2=e$. It follows that $q^2=q$ and $eq=qe=0$.
The augmentation of $e$ is
$\left(\begin{smallmatrix}4&-12\\1&-3\end{smallmatrix}\right)$.
Each of $e,q$ has exactly 42 signed group-labelled matrix units
when coefficients are expanded with their multiplicities. The full
coefficient arrays and the four projector identities are checked
independently by \path{code/d_c23_short_lift.py} and by the final
table constructor.

We verify that changing the augmentation lift preserves the required
nonvanishing conclusion. Under $R\to\mathcal O_L=\Z[\zeta_{23}]$,
the norm element $N$ vanishes. The image of $e$ is therefore exactly
\[
 \bar e=\begin{pmatrix}4&\omega-1\\2\omega&-3\end{pmatrix},
\]
the same cyclotomic projector used in
Proposition~\ref{D:C23-projector}. Let
$K=\Q(\omega)=\Q(\sqrt{-23})$ and
$I=(2,\omega)\subset\mathcal O_K=\Z[\omega]$.
The ideal has norm two and is not principal: the norm form
\[
 a^2+ab+6b^2=(a+b/2)^2+23b^2/4
\]
does not take the value two on integers. The element $\omega-2$
has norm eight and is absent from the conjugate prime
$(2,\omega-1)$, so $I^3=(\omega-2)$. Thus $[I]$ has exact order
three. Since $[L:K]=11$, taking ideal norm shows that
$I\mathcal O_L$ cannot be principal: otherwise $I^{11}$ would
be principal, contrary to its order.
Writing
\[
 p=\binom{2}{\omega},\qquad
 t=\bigl(2,(\omega-1)/2\bigr)
\]
gives $tp=1$ and $pt=\bar e$. Consequently
$\operatorname{im}\bar e\cong(I\mathcal O_L)^{-1}$.
Its determinant line is nontrivial, so it is not stably free.
Extension of scalars would send a zero reduced projective class
to zero. Therefore
\[
 [\operatorname{im}e]\ne0\quad\text{in }\widetilde K_0(R).
\]
This argument asserts nonvanishing for the new lift; equality with
the old integral projective class is not needed.

\paragraph{The finite base and all multiplication disks.}
Let $L$ be the join of two 23-gon circles. Number the first circle
by $0,\ldots,22$ and the second by $23,\ldots,45$, with their
cyclic edge orders. The group $C_{23}$ acts by simultaneous
addition modulo 23 in both layers. Since a nonidentity rotation
fixes no point in either circle, this action on
$|L|=S^1*S^1\cong S^3$ is free. For each nonempty face orbit
choose its lexicographically least face. Define $K_{23}$ by these
representatives and the arrows
\[
 (F,H,t):F\longrightarrow H\quad(t+F\subseteq H),
 \qquad
 (H,J,u)(F,H,t)=(F,J,u+t),
\]
with labels in $\Z/23$. There are 96 objects and 744 arrows;
its nondegenerate nerve counts are $(96,648,1104,552)$.
The face poset of $L$ is its free covering category. Hence
$BK_{23}\cong|\operatorname{sd}L|/C_{23}$ has fundamental group
$C_{23}$ and universal cover $S^3$.

For a label $a\in\{1,\ldots,22\}$ take the path
$(0,1,\ldots,a)$ in the first polygon. The identity label uses
the padded constant path $(0,0)$. For an inner label $h$ and
outer label $g$, concatenate the first path with its translate
$h+\operatorname{path}(g)$, then close it by the reverse of the
chosen path for $(h+g)\bmod23$. Cone this parameter polygon to
vertex 23 in the second layer. Every labelled triangle is a face
of the join, possibly with a collapsed boundary edge. This fixes
all 529 multiplication disks; their boundary lengths lie between
3 and 65. For an inclusion of parameter faces $F\subseteq H$,
choose the unique transport labels $t_F,t_H$ from their normalized
image faces and use the arrow label $t_F-t_H$. This supplies exact
functors to $K_{23}$ and exact boundary equalities. Their complete
native tables are saved in
\path{code/outputs/d_c23_coefficient_context_certificate.json.gz}.

Use $Y=NK_{23}\vee A_1^{\wedge5}\vee A_1^{\wedge5}$ and the same
signed elementary cospans, flat pinches, permutation grids, and
simultaneous cancellation cylinders as in
Theorem~\ref{D:Q8-final-table}. Expanding each of $eq$ and $qe$
gives 1764 terms. The deterministic pairing order is lexicographic
in the target index, group label, and source index, with occurrences
paired in their original order. For $eq$ this yields 405 opposite
pairs and 954 zero products; for $qe$ it yields 416 opposite pairs
and 932 zero products. Each specified nullhomotopy has 50 cospan
roofs and 49 actual comparison maps. All occurrences of a group
product use the corresponding fixed fan disk.

\paragraph{The arithmetic table and its realization.}
The finite simplicial normal forms and ordinal pullback rules are
unchanged. For $A_n$ use this fixed arithmetic order in degree $k$:
the basepoint, the $n-1$ constant interior vertices, the $n$ constant
edges, then the $2nk$ crossing strings ordered by edge, endpoint
side, and the number of initial vertices. Thus all ranks and
unranks are given by integer arithmetic. The parameter interval
$J_n$ lists first its vertices, then its edges; its arrows list
identities, left incidences, then right incidences. The optional
backend \path{code/d_arithmetic_intervals.py} implements precisely
these orders and the same composition functions. The displayed
constructor is run in a fresh Python interpreter, installing this
backend before constructing the new instance; its process-local
factory replacements are not mixed with a quaternion-table run.

Apply the finite simplex-category replacement $\mathcal O_8$ to
each coefficient space, then the ordinal, product, span, and
height-one constructors already specified in
Theorem~\ref{D:Q8-final-table}. Every coefficient roof has dimension
at most seven; the final relative nullhomotopy cylinders have
dimension at most eight. Thus the proved truncation equivalence
$B\mathcal O_8X\simeq|X|$ applies to every input. The same two
parity monoids, suspension comparison, tensor-zero contraction,
base-path collars, and intrinsic pinch produce the final finite
category. Its 3897 homogeneous type blocks have exactly
\[
\begin{split}
 \#\{\text{typed pairs}\}
  &=288\,986\,603\,572\,868\,508\,990\,808\,637\,052,\\
 \#\{\text{typed triples}\}
  &=109\,162\,136\,788\,728\,112\,368\,563\,487\,083\,799\,108.
\end{split}
\]
Every lookup is a finite recursion on these constructor types;
there is no subdivision search or unstated choice in the table.

\paragraph{Chains, domination, and obstruction.}
All maps are over $BK_{23}$, with their under-base identifications
retained as the explicit natural-transformation paths. Van Kampen
therefore gives $\pi_1(B\mathcal W_{23})=C_{23}$. The same parity
calculation gives relative chains $R^2$ in each degree at least
seven, and alternates $e,q$. After the first cone and its explicit
deformation retraction, the second attaching map has coordinates
\[
 x\longmapsto\bigl(\alpha_n(1-f_n)x,\,\beta_nf_nx\bigr),
 \qquad \alpha_n,\beta_n\in\{1,-1\}.
\]
Its inverse on $(1-f_n)R^2\oplus f_nR^2$ is
$(u,v)\mapsto\alpha_nu+\beta_nv$. Thus every positive relative
degree contracts, leaving $qR^2$ in degree seven. The coefficient
nullhomotopies have zero degree-five-to-six cellular homotopy
term, and the extra tensor-zero coherence factors through the
zero relative coefficient quotient, exactly as in the first
table. No additional matrix commutativity or sign assumption
enters the cone calculation.
Consequently
\[
 C_*(\widetilde{B\mathcal W_{23}})\simeq
 C_*(\widetilde{BK_{23}})\oplus qR^2[7],
\]
where the base complex has free ranks $(96,648,1104,552)$ in
degrees zero through three.

The finite seven-skeleton of the nerve dominates the whole space:
its inclusion is seven-connected, and all section-obstruction
groups in degrees at least eight vanish for every local coefficient
system by this bounded projective chain model. A mapping telescope
of the resulting homotopy idempotent on that finite skeleton gives
a CW model of dimension at most eight. Finally Wall's obstruction
is
\[
 w(B\mathcal W_{23})=-[qR^2]=[eR^2]\ne0.
\]
Wall's theorem therefore excludes a finite CW model.

\paragraph{Executed checks.}
All 529 native fan disks were checked, with 911\,881 composable
triples and 107\,272 exact boundary equations. The independent
elementary-coefficient checker covered all 4232 group/sign/index
cases and 338\,560 endpoint equations. Each complete nullhomotopy
passed 2940 selected endpoint/pullback equations. The final run
passed 558 base-retraction equations, 468 naturality equations,
and 15\,588 identity/associativity equations at the ends of all
3897 type blocks. The last checks sample the large table; the
universal constructor proofs establish its laws for every index.
The exact record is
\path{code/outputs/d_c23_wall_category_summary.json}.
\end{proof}


\begin{corollary}[The specified cases in D2]\label{D:D2-outcomes}
Finite-dimensional classifying-space type implies finite CW type for
finite EI categories and for finite categories with fundamental group
$C_2$. The assertion fails for each of $Q_8$ and $C_{23}$.
\end{corollary}
\begin{proof}
The positive statements are Theorems~\ref{D:EI-finite}
and~\ref{D:C2-finite}. For the negative statements use the executed
indexed tables of Theorems~\ref{D:Q8-final-table}
and~\ref{D:C23-final-table}; each has its nonzero Wall class
computed explicitly.
\end{proof}

The induced syzygy of Proposition~\ref{D:syzygy} consequently need
not be stably free. The four-arrow example explains why ordinary
category-module projectivity could not establish that conclusion;
the executed counterexample shows that the proposed induced
stable-freeness conclusion itself is false.

\subsection{Two objects can carry nontrivial deck coupling}

\begin{lemma}[Representing a corner map by a finite biset]
\label{D:corner-biset}
Let $P,M$ be finite monoids and let $\phi:P\to eMe$ be a unital
homomorphism, where $e\in M$ is idempotent. Regard $Me$ as an $(M,P)$-biset
using $\phi$ on the right. Its category of elements has classifying space
homotopy equivalent to $BP$; under this equivalence its two projections
induce the identity of $BP$ and the corner map $BP\to BM$ defined by $\phi$.
Consequently, for any finite family $\phi_i:P\to e_iMe_i$, the two-object
category with endomorphism monoids $P,M$, no arrows in the reverse
direction, and cross-biset $\coprod_i Me_i$ realizes
\[
 \operatorname{hocolim}\left(
 BP\ \longleftarrow\ \coprod_i BP\ \longrightarrow\ BM\right),
\]
with left maps the identity and right maps the prescribed corner maps.
\end{lemma}
\begin{proof}
For one component, the category of elements has objects $x\in Me$ and
arrows $x\to x'$ given by pairs $(p,m)\in P\times M$ satisfying
$x'\phi(p)=mx$. Projection to $P$ is a Grothendieck fibration:
the cartesian lift over $p$ with target $x'$ has source $x'\phi(p)$
and second component $1$. Its fiber is the action category $M\ltimes Me$.
The $M$-set $Me$ is a retract of the regular $M$-set $M$, by
$m\mapsto me$ and inclusion. The action category $M\ltimes M$ has initial
object $1$. Therefore $B(M\ltimes Me)$ is a retract of a contractible
space and is contractible. The fiber-to-comma adjunction for the displayed
fibration, or Quillen's Theorem A applied to these cartesian lifts, makes
projection to $BP$ a homotopy equivalence.

In the Karoubi category of $M$, the arrows $x:e\to1$ give a natural
transformation between the two projected functors: its naturality
identity is exactly $x'\phi(p)=mx$. This identifies the other projection
with the specified corner map. The last assertion is the collage
homotopy-pushout lemma checked in Lemma~\ref{lem:B-prelim}.
\end{proof}

\begin{proposition}[A two-object degree-two mapping torus]
\label{D:degree-two-torus}
The mapping torus of a degree-two self-map of $S^2$ is the classifying
space of a finite category with exactly two objects and aperiodic
endomorphism monoids. Its fundamental group is $\mathbb Z$, and
\[
 H_2(\widetilde X;\mathbb Z)\cong\mathbb Z[1/2].
\]
In particular, finite two-object categories are not restricted to wedges
of circles and simply connected summands, nor to finite-order monodromy.
\end{proposition}
\begin{proof}
Let $P$ be the five-element rectangular-band monoid realizing $S^2$.
Let $L$ be the monoid obtained by adjoining an identity to the
$3\times2$ rectangular band. Then
$BL\simeq S^2\vee S^2$ by Proposition~\ref{prop:B-band-wedge}.
The two unital embeddings $P\to L$ selecting row pairs $\{0,1\}$ and
$\{0,2\}$ induce the two integral Hurewicz generators, denoted $a,b$.
Apply the checked finite mapping-cone construction to a sphere map
representing $2a-b$. It gives a finite aperiodic monoid $M$ and a full
corner $eMe\cong L$, with
\[
 BM\simeq (S^2\vee S^2)\cup_{2a-b}D^3\simeq S^2,
\]
compatibly with the map from $BL$.
The last equivalence follows by a change of integral Hurewicz basis on
the wedge: the primitive vector $2a-b$ is one basis vector and its
attaching cell cancels that sphere. The images of $a,b$ in the remaining
$H_2\cong\mathbb Z$ are respectively one and two times a generator.

Compose the two embeddings with $L\cong eMe$. Lemma~\ref{D:corner-biset}
applied to $Me\amalg Me$ gives a finite two-object category whose
classifying space is the homotopy coequalizer of maps
$S^2\rightrightarrows S^2$ of degrees one and two. Replacing the
degree-one map by a homotopy equivalence identifies this with the
degree-two mapping torus. Its universal cover is the bi-infinite mapping
telescope of degree-two maps, so its second homology is the direct limit
of multiplication by two on $\mathbb Z$, namely $\mathbb Z[1/2]$.

Already the product of the two-parallel-arrow category with $P$ gives
a simpler two-object model of $S^1\times S^2$. The degree-two example
shows why eventual periodicity of an endomorphism of a fixed finite
monoid does not restrict all two-object realizations: the two maps here
have different source and target monoids.
\end{proof}

\subsection{Simultaneous finite realization of maps by full corners}

\begin{lemma}[A wedge construction preserving full corners]
\label{D:corner-wedge}
Let $P,Q$ be finite aperiodic monoids whose classifying spaces have finite
homotopy type. There is a finite aperiodic monoid $R$ and full-corner
embeddings of $P,Q$ into $R$ such that
\[
 BR\simeq BP\vee BQ
\]
under the two corner maps. Each left and right ideal joining a corner to
$1_R$ is a finite projective act over the corresponding corner monoid.
Here a finite projective act means a retract of a finite free act, or
equivalently a finite disjoint union of idempotent-generated principal
ideals.
\end{lemma}
\begin{proof}
The trivial factor can be omitted, retaining a minimal-ideal idempotent
for its one-element corner. Otherwise let
$K_P=I_P\times J_P$, $K_Q=I_Q\times J_Q$ be the minimal rectangular
bands, and choose $e_P\in K_P$, $e_Q\in K_Q$. Inside $P\times Q$ put
\[
 T=\{(1,1)\}\cup(P\times\{e_Q\})\cup(\{e_P\}\times Q)
       \cup(K_P\times K_Q).
\]
This is a submonoid: a product crossing the middle two pieces lies in
$K_P\times K_Q$. Its minimum is that rectangular band. The idempotents
$f=(1,e_Q)$ and $g=(e_P,1)$ satisfy
$fTf\cong P$, $gTg\cong Q$, because
$e_PK_Pe_P=\{e_P\}$ and $e_QK_Qe_Q=\{e_Q\}$.

Write $A_P=\mathbb ZP/\mathbb ZK_P$ and define $A_Q$ similarly. In
$A_T=\mathbb ZT/\mathbb Z(K_P\times K_Q)$ the idempotents
$f,g,h=1-f-g$ are central and orthogonal, giving
\[
 A_T\cong A_P\times A_Q\times\mathbb Z.
\]
For the augmentation modules of the row sets, the three summands are
\[
 fV_T\cong V_P,\qquad gV_T\cong V_Q,\qquad
 hV_T\cong V_P\otimes_{\mathbb Z}V_Q;
\]
the column augmentation modules satisfy the analogous identities.
These statements follow by writing the free row module as
$(\mathbb Z\oplus V_P)\otimes(\mathbb Z\oplus V_Q)$, using the chosen
rows as splittings of augmentation. The actions of $f,g$ project to the
first and second augmentation summands, while $h$ projects to their tensor
product.

Lemma~\ref{lem:B-band-Tor} and the product decomposition of $A_T$ now give
\[
 H_n(BT;\mathbb Z)\cong H_n(BP;\mathbb Z)\oplus H_n(BQ;\mathbb Z)
       \quad(n\ge3),
\]
and the same formula in degree two with an additional free summand of rank
\[
 d=(|I_P|-1)(|J_P|-1)(|I_Q|-1)(|J_Q|-1).
\]
The corner inclusions are split by the two coordinate projections
$T\to P,Q$; each opposite composite is constant at a minimal idempotent.
Thus their images form a split direct sum in homology. Since all groups
in each degree are finitely generated, the displayed abstract computation
shows that its complement is zero outside degree two and is
$\mathbb Z^d$ in degree two. All three classifying spaces are simply
connected, since their minimum ideals are bands. Choose Hurewicz spheres
for a basis of this complement. The resulting map
\[
 BP\vee BQ\vee\bigvee^d S^2\longrightarrow BT
\]
is an integral homology isomorphism between simply connected CW complexes,
hence a homotopy equivalence.

If $d>0$, use the finite mapping-cone construction to kill precisely these
$d$ sphere summands. It yields the required $R$ under $T$; if $d=0$, take
$R=T$. The embeddings remain full corners under this construction.

For completeness, before the cone attachment the right $P$-act $Tf$ is
\[
 P\ \amalg\!\coprod_{q\in Qe_Q\setminus\{e_Q\}} K_P.
\]
Every right $\mathcal R$-class of $K_P$ is an idempotent-generated right
ideal of $P$, so this act is projective. The analogous description proves
projectivity of $fT$, and likewise for $g$. In the explicit monoid
$M(Y,u)$ of Part B, the new left and right ideals at the old corner are
free acts over the old monoid: they are respectively $L\times[n]$ and
$Y_n\times L$. Iterating cone attachments consequently preserves
projectivity of all the older side acts, by the identities
$R f\cong R e\otimes_{eRe}(eRe)f$ and its left-hand counterpart.
This proves all assertions.
\end{proof}

\begin{lemma}[Relative cell trading]
\label{D:relative-cell-trading}
For a map $f:T\to X$ between finite connected simply connected CW
complexes, there is a finite relative CW complex $Z$ built from $T$ using
only cells of dimension at least two, and a homotopy equivalence $Z\simeq X$
whose restriction to $T$ represents $f$.
\end{lemma}
\begin{proof}
Replace $f$ by a cellular map and start with its finite mapping cylinder.
A relative maximal tree removes the relative zero-cells and their
connecting one-cells. Each remaining relative one-cell, completed to a
loop using paths in $T$, is null-homotopic in the mapping cylinder because
the latter has the homotopy type of $X$.

For every such one-cell add a two-cell along this completed loop and a
three-cell along the sphere formed by the new two-cell and a chosen null
disk in the old complex. These finitely many two/three-cell pairs leave
the homotopy type unchanged under $T$: the new two-cell adds a sphere,
and the new three-cell kills that sphere with coefficient one. Cellular
approximation places the chosen null disks in the old two-skeleton.
The subcomplex consisting of $T$, the relative one-cells, and the new
two-cells deformation retracts to $T$, since each new two-cell runs once
over its designated one-cell and otherwise lies in $T$. Replace this
subcomplex by $T$, transporting the remaining attaching maps along the
retraction. Homotopy invariance of pushouts along CW cofibrations preserves
the homotopy type under $T$. Only the old cells of dimension at least two
and the new three-cells remain.
\end{proof}

\begin{proposition}[Simultaneous full-corner realization]
\label{D:simultaneous-corners}
Let $T_1,\ldots,T_k,X$ be finite connected simply connected CW complexes,
and let $f_i:T_i\to X$ be a finite family of maps. Fix finite aperiodic
monoids $P_i$ realizing the $T_i$. There is a finite aperiodic monoid $M$
with full corners $e_iMe_i\cong P_i$ such that $BM\simeq X$ and the corner
maps represent the $f_i$. The side acts $Me_i$ and $e_iM$ are finite
projective $P_i$-acts.
\end{proposition}
\begin{proof}
Iterate Lemma~\ref{D:corner-wedge} to obtain a finite aperiodic monoid $P$
realizing $\bigvee_i T_i$, under full corners $P_i$, with projective side
acts. The maps $f_i$ assemble to a map from this wedge to $X$; basepoints
can be chosen compatibly because $X$ is simply connected. Apply
Lemma~\ref{D:relative-cell-trading}. Its finite sequence of cells of
dimension at least two can be attached successively using the checked
construction in Proposition~\ref{prop:B-finite-realization}: each attaching
sphere is connected, and each current monoid is aperiodic. Compatibility
under the old corner preserves the specified maps. At each step the old
corner is full and the new side acts are free over the old monoid, as
checked in the proof of Lemma~\ref{D:corner-wedge}. Composition therefore
preserves fullness and projectivity of the original $P_i$ corners.
\end{proof}

\begin{lemma}[A finite biset for a correspondence]
\label{D:projective-correspondence}
Suppose $P$ embeds as full corners $e_A M_A e_A$ and $e_B M_B e_B$, with
finite projective side acts. The finite $(M_B,M_A)$-biset
\[
 S=M_Be_B\otimes_P e_A M_A
\]
realizes the correspondence $BM_A\leftarrow BP\to BM_B$: its two-object
collage has the homotopy type of this homotopy pushout, compatibly with
the two vertex maps. A disjoint union of such bisets realizes the
corresponding disjoint union of edge spaces.
\end{lemma}
\begin{proof}
In the nerve of a two-object collage, a simplex crossing from the first
object to the second has a string of $M_A$ arrows, one element of $S$,
and a string of $M_B$ arrows. Splitting its barycentric coordinates at
the cross arrow gives the double mapping cylinder whose edge space is
the realization of this two-sided double bar construction.

The augmentation
\[
 B(M_Be_B,P,e_A M_A)\longrightarrow
 M_Be_B\otimes_P e_A M_A
\]
is a homotopy equivalence. Indeed it is so when its first right $P$-act
is free, by the extra degeneracy of the usual bar construction, and
finite projective acts are retracts of finite free acts. This equivalence
respects the outer actions. Substitute this bar resolution for $S$ in
the edge space and reassociate the bar directions. The two outer factors
become the action categories of the representable acts $M_Be_B$ and
$e_A M_A$; their classifying spaces are contractible, by the retract
argument of Lemma~\ref{D:corner-biset} and its dual. Replacing these
contractible factors by points leaves $B(*,P,*)=BP$. The two projections
are the corner maps, as follows from the same natural transformations
in the Karoubi categories used in Lemma~\ref{D:corner-biset}.
Realization and the bar reassociation preserve these comparisons because
all the simplicial spaces involved are CW spaces with cofibrant
degeneracy inclusions. Disjoint unions are preserved throughout.
\end{proof}

\begin{theorem}[Two objects realize finite graphs of simply connected spaces]
\label{D:graphs-two-objects}
Every finite connected graph of finite connected simply connected CW
spaces, with arbitrary continuous edge attaching maps, has the homotopy
type of the classifying space of a finite category with exactly two
objects and aperiodic endomorphism monoids.
\end{theorem}
\begin{proof}
Take the homotopy colimit over a maximal tree of the underlying graph.
It is a finite simply connected CW complex $Y$. The remaining edges
give finitely many pairs of maps $T_i\rightrightarrows Y$, with each
$T_i$ finite and simply connected. The original graph of spaces is the
homotopy pushout of two copies of $Y$ along the disjoint union
\[
 Y\amalg T_1\amalg\cdots\amalg T_r,
\]
where the first edge maps identically to both copies of $Y$, and the
other edge maps are the specified pairs. Collapsing the identity-edge
mapping cylinder recovers the original model.

Choose a finite aperiodic monoid $P_0$ for $Y$ and monoids $P_i$ for the
$T_i$. Apply Proposition~\ref{D:simultaneous-corners} separately to the
two families of maps into $Y$, including the identity map from $Y$.
This gives $M_A,M_B$ with the required common source corners and
projective side acts. Take the disjoint union of the finite bisets in
Lemma~\ref{D:projective-correspondence}, one for each of these $r+1$
edges. Their two-object collage realizes the required homotopy pushout.
\end{proof}

\begin{lemma}[A graph as the one-skeleton]
\label{D:rose-skeleton}
If $X$ is a finite connected CW complex with fundamental group $F_r$,
then $X$ has a finite CW model whose two-skeleton is
$R_r\vee\bigvee_{j=1}^a S^2$, where $R_r$ is a rose with $r$ petals.
\end{lemma}
\begin{proof}
Choose a map $R_r\to X$ inducing the given isomorphism on fundamental
groups, and take its finite mapping cylinder. The proof of
Lemma~\ref{D:relative-cell-trading} applies whenever the map on fundamental
groups is surjective: complete a relative one-cell to a loop, choose a
loop in the source representing the same element of the target fundamental
group, and use their difference as the attaching loop of the new two-cell.
This difference is null-homotopic in the target, so the compensating
three-cell exists. The new two-cell still traverses its designated
relative one-cell exactly once. Thus the same finite cell trading produces
a finite relative CW model over $R_r$ with no relative zero- or one-cells.

Each remaining two-cell attaches along a loop in $R_r$ that is null in
the total complex. Injectivity of $\pi_1(R_r)\to\pi_1(X)$ makes that loop
null-homotopic already in $R_r$. Replacing these finitely many attaching
maps by constant maps, and transporting the higher attaching maps through
the resulting equivalence of two-skeleta, gives the asserted model.
\end{proof}

\begin{lemma}[Enlarging pieces over a tree]
\label{D:tree-enlargement}
Let $Z$ be a finite graph of finite connected simply connected CW spaces
over the rose $R_r$, and let $n\geq3$. For any map
$\alpha:S^{n-1}\to Z$, the space $Z\cup_\alpha D^n$ again has a model of
this form over the same rose.
\end{lemma}
\begin{proof}
We may assume $r\geq1$ and that all attaching maps are cellular. Write
$G=F_r$ and let $T$ be the Cayley tree for its displayed free basis.
The universal cover $Y$ of $Z$ is the tree of the prescribed simply
connected vertex and edge spaces. It has its usual projection $q:Y\to T$:
vertex spaces map to vertices, and edge mapping cylinders map to edges.
For every nonempty finite subtree $K\subset T$, the inverse image
$q^{-1}(K)$ is a finite simply connected CW complex. This follows by
successively adjoining leaf vertex spaces along connected simply
connected edge spaces, using van Kampen; it also follows directly from
the graph-of-spaces homotopy colimit.

Lift $\alpha$ to $\widetilde\alpha:S^{n-1}\to Y$, and choose a finite
subtree $S\subset T$ containing $q\widetilde\alpha(S^{n-1})$. Compactness
and local finiteness of $T$ permit this choice. For an integer $R\geq1$
large enough that some combinatorial radius-$R$ ball contains $S$, put
\[
 D_g=q^{-1}(B_R(g)) \qquad(g\in G).
\]
Here $B_R(g)$ is the full subtree on vertices at distance at most $R$
from $g$. For every edge $[g,h]$ of $T$, put
$E_{gh}=D_g\cap D_h=q^{-1}(B_R(g)\cap B_R(h))$.
Both types of pieces are finite and simply connected, because the
intersection of these two balls is a nonempty finite subtree.

The tree homotopy colimit with vertex pieces $D_g$ and edge pieces
$E_{gh}$ maps naturally to $Y$, and this map is a homotopy equivalence.
Here is a cellular justification that also keeps track of the finite
quotient. For each open cell $\sigma$ of the original tree of spaces,
its set of membership vertices
\[
 A_\sigma=\{g:\sigma\subset D_g\}
\]
is a nonempty finite subtree of the indexing tree $T$. Indeed $q(\sigma)$
is either a vertex or an open edge, and the membership set is respectively
one ball of centers or the intersection of the two balls of centers
determined by its endpoints. At the cellular stage for $\sigma$, the
homotopy colimit therefore adjoins
$D^{\dim\sigma}\times |A_\sigma|$ along
$S^{\dim\sigma-1}\times |A_\sigma|$; the projection to the corresponding
cell of $Y$ is a homotopy equivalence of these attaching pairs because
$|A_\sigma|$ is contractible. Induction over the skeleta and the CW gluing
lemma prove the assertion. The construction is $G$-equivariant, with no
cell stabilizers. The same cellular argument on cell orbits proves that
the quotient graph homotopy colimit is homotopy equivalent to $Z$.

For each translated attaching map $\gamma\widetilde\alpha$, define
\[
 A_\gamma=\{g:\gamma S\subset B_R(g)\}
          =\bigcap_{v\in\operatorname{Vert}(\gamma S)}B_R(v).
\]
This is a nonempty finite subtree: it is nonempty by the choice of $R$
and translation, and intersections of convex subtrees of a tree are
convex. Attach a copy of its $n$-cell to every $D_g$ for $g\in A_\gamma$,
and attach the same cell to $E_{gh}$ whenever $[g,h]\subset A_\gamma$.
The boundary map lies in all the required pieces by the definition of
$A_\gamma$. Call the enlarged vertex and edge pieces $D'_g,E'_{gh}$.
They are simply connected since only cells of dimension at least three
have been added. They are finite: for a fixed $g$, only finitely many
translates $\gamma S$ lie in $B_R(g)$, since the action on vertices is
free and the ball is finite.

The new tree homotopy colimit is obtained from the old one by adjoining,
for each $\gamma$, the pair
\[
 (D^n\times |A_\gamma|,
  S^{n-1}\times |A_\gamma|).
\]
Its natural map to $Y$ with all translates of the new cell attached
collapses the contractible tree $|A_\gamma|$. Homotopy invariance of the
attaching pushouts, or the same cellular argument, makes this map a
homotopy equivalence. Everything is $G$-equivariant. Taking the quotient
gives a finite graph of the simply connected spaces $D'_1$ and
$E'_{1,s}$, one edge for each free generator $s$, and a homotopy
equivalence from its total space to $Z\cup_\alpha D^n$.
\end{proof}

\begin{theorem}[Finite complexes with free fundamental group]
\label{D:free-pi-two-objects}
A finite connected CW complex $X$ is homotopy equivalent to the
classifying space of a finite category with exactly two objects if and
only if its fundamental group is free. In the sufficient direction the
endomorphism monoids can both be chosen aperiodic.
\end{theorem}
\begin{proof}
For sufficiency, Lemma~\ref{D:rose-skeleton} starts a finite CW model of
$X$ with $R_r\vee\bigvee_a S^2$. This is a graph of simply connected
finite spaces over $R_r$: use $\bigvee_a S^2$ as the vertex space and a
point for each edge space. Attach all the finitely many higher cells,
one at a time, using Lemma~\ref{D:tree-enlargement}. The resulting finite
graph of finite simply connected spaces is realized by the two-object
category of Theorem~\ref{D:graphs-two-objects}.

For necessity, Theorem~\ref{A:bounded-categories} kills every endomorphism
loop in the fundamental groupoid of any finite category with finite
CW type. If both cross-hom sets of a two-object category are nonempty,
choose arrows $a:0\to1$ and $b:1\to0$. Their composites are endomorphisms,
so $b=a^{-1}$ in the fundamental groupoid; any other cross arrow equals
$a$ or $b$ by the same argument. The fundamental group is then trivial.
Otherwise the connected category is triangular. After its endomorphisms
have been killed, its fundamental-group presentation has one generator
for each equivalence class of cross arrows under left and right
endomorphism actions and no other relations. Choosing one such class
as a spanning-tree edge leaves a free group on the remaining classes.
\end{proof}

\begin{corollary}[The full two-object classification]
\label{D:two-object-remaining}
The connected finite-dimensional CW homotopy types realized by finite
categories with exactly two objects are precisely the finite connected
CW homotopy types with free fundamental group.
\end{corollary}
\begin{proof}
Theorem~\ref{D:two-objects} proves that a finite two-object category with
finite-dimensional classifying-space type has finite CW type and free
fundamental group: finite domination, vanishing of
$\widetilde K_0(\mathbb ZF_r)$, and Wall's theorem supply finiteness.
The converse is Theorem~\ref{D:free-pi-two-objects}.
\end{proof}

\begin{remark}\label{D:tree-repair-scope}
The tree enlargement proof uses no homology-preserving plus construction
on the individual vertex or edge pieces. Their homology is allowed to
change; the product with a contractible membership subtree controls the
homotopy type of the total space.
\end{remark}

\begin{corollary}\label{D:two-object-disconnected}
The disconnected finite-dimensional types realized by categories with
exactly two objects are precisely disjoint unions of two nonempty
finite connected simply connected CW homotopy types.
\end{corollary}
\begin{proof}
A disconnected two-object category has no cross arrows and is the
disjoint union of its two endomorphism monoids. Section~\ref{sec:A}
forces each finite-dimensional component to be simply connected and
of finite CW type. Conversely Theorem C, checked in Section~\ref{sec:B},
realizes the two given simply connected finite types separately by
finite monoids; their disjoint union is the required category.
\end{proof}

